\documentclass[11pt]{article}

\usepackage[T1]{fontenc}
\usepackage[margin=1in]{geometry}
\usepackage{amsmath,amssymb,mathtools,amsthm}
\usepackage{booktabs,tabularx,array}
\IfFileExists{lmodern.sty}
  {\usepackage{lmodern}\usepackage{microtype}}
  {\usepackage[expansion=false]{microtype}}
\usepackage{enumitem}
\usepackage[hidelinks]{hyperref}

\setlength{\parindent}{0pt}
\setlength{\parskip}{0.65em}
\setlist[itemize]{leftmargin=*,itemsep=0.2em,topsep=0.3em}
\setlist[enumerate]{leftmargin=*,itemsep=0.3em,topsep=0.3em}
\renewcommand{\arraystretch}{1.18}
\setcounter{tocdepth}{1}

\newcommand{\AP}{\operatorname{AP}}
\newcommand{\CNAP}{\operatorname{CNAP}}
\newcommand{\Regime}{\mathcal{R}}
\newcommand{\PSet}{\Pi_\star}
\newcommand{\E}{\mathbb{E}}
\newcommand{\Prb}{\mathbb{P}}
\newcommand{\iid}{\mathrel{\overset{\mathrm{iid}}{\sim}}}
\newcommand{\RobCNAP}{\underline{\CNAP}_{\PSet}}
\newcommand{\RobDelta}{\underline{\Delta}_{\PSet}}
\newcommand{\SRob}{S_{K,\PSet}^{\mathrm{AP}}(\Regime)}
\newcommand{\SRobHat}{\widehat S_{K,\PSet}^{\,\mathrm{AP}}}
\newcommand{\SRobNullHat}{\widehat S_{K,\PSet}^{\,\mathrm{AP,null}}}
\newcommand{\SRobCalHat}{\widehat S_{K,\PSet}^{\,\mathrm{AP,cal}}}
\newcommand{\APta}{\widetilde{\AP}}
\newcommand{\CNAPta}{\widetilde{\CNAP}}
\newcommand{\Dta}{\widetilde{\Delta}}
\newcommand{\RobDta}{\underline{\widetilde{\Delta}}_{\PSet}}
\newcommand{\Rank}{\mathsf{O}}
\newcommand{\RefSet}{\mathcal{F}}

\newtheorem{definition}{Definition}
\newtheorem{proposition}{Proposition}
\theoremstyle{remark}
\newtheorem*{remark}{Remark}

\title{Supported Predictive Performance Across Target Prevalences\\
\large A Replication-Based Method for Average Precision}
\author{Marcus A. Thomas}
\date{Methodological proposal}

\begin{document}

\maketitle

\begin{abstract}
Predictive performance serves two tasks: interpreting the evidence for a model considered on its own and informing a decision between models evaluated together. Average precision makes both tasks sensitive to the prevalence of the population receiving the model. Prior standardization evaluates precision at each prevalence in a declared target set $\PSet$, and chance normalization creates a common CNAP scale with population random ranking at zero and perfect ranking at one. Each evaluation is summarized by its lowest CNAP over $\PSet$. For $K$ independent executions of an evaluation-generating regime $\Regime$, supported performance is
\[
S_{K,\PSet}^{\mathrm{AP}}(\Regime)
=
\E_{\Regime}
\left[
\min_{1\leq j\leq K}
\inf_{\pi\in\PSet}\CNAP_\pi(D_j)
\right].
\]
The level must survive every declared prevalence and every replication, and widening $\PSet$ or raising $K$ gives the claim more opportunities for defeat. A stratified bootstrap estimates the quantity. For an isolated model, conditional permutation locates it against the chance level of its own finite design, and the calibrated ratio is what the report carries. For two models on the same units, the operator is applied to the lowest paired CNAP difference over $\PSet$, computed under a convention that averages over the orderings within each tie block. That convention makes the differential design anchor vanish identically wherever the outcome labels of one evaluation are equally likely in every arrangement, so the comparison keeps the raw scale and estimates no correction, at the cost of being conservative. A replacement claim requires the supported advantage to exceed a declared margin, with a proposed outer lower bootstrap bound above that margin as the inferential criterion.
\end{abstract}

\begingroup
\small
\setlength{\parskip}{0pt}
\tableofcontents
\endgroup

\begin{table}[ht]
\centering
\caption{Principal symbols.}
\label{tab:notation}
\begin{tabularx}{\textwidth}{@{}>{$}l<{$}X@{}}
\toprule
\text{Symbol} & Meaning \\
\midrule
\pi & A target-population prevalence \\
\PSet & The declared set of target prevalences that may challenge the claim \\
\Regime & The evaluation-generating regime, declaring what varies across replications \\
K & The number of independent replications that must retain the claim \\
\AP_\pi(D) & Prior-standardized average precision of evaluation $D$ at prevalence $\pi$ \\
\CNAP_\pi(D) & Chance-normalized $\AP_\pi$, zero at population random ranking and one at perfect ranking \\
\RobCNAP(D) & Lowest CNAP over $\PSet$, the prevalence-robust effect of one evaluation \\
\pi^\dagger(D) & A prevalence attaining that lowest value \\
Z_{\PSet} & The prevalence-robust effect of one replication drawn from $\Regime$ \\
S_K(T;\Regime) & Replication-survival operator, the mean of $\min_{j\leq K}T_j$ \\
\SRob & Prevalence-robust supported CNAP, the primary population quantity \\
\SRobHat & Its estimate from the bootstrap replication array \\
\SRobNullHat & The robust conditional chance level from permuted outcomes \\
\SRobCalHat & The calibrated value, reported for an isolated model \\
C_{\mathrm{prev}},\ C_{\mathrm{rep}} & The prevalence-challenge and replication-disagreement costs \\
\Delta_\pi(D) & Paired CNAP difference between models $A$ and $B$ at prevalence $\pi$ \\
\RobDelta(D) & Lowest paired difference over $\PSet$ \\
\CNAPta_\pi(D) & Tie-averaged CNAP, the convention required for paired comparison \\
\RobDta(D) & Lowest tie-averaged paired difference over $\PSet$ \\
d & The smallest advantage that would justify replacing one model with another \\
\bottomrule
\end{tabularx}
\end{table}

\newpage

\section{Population conditions, corroboration, and model decisions}
\label{sec:motivation}

Many prediction systems rank cases so that a selected set can receive further attention. Candidate peptides may be selected for a vaccine, radiological images may be assigned for review, and customers may be prioritized for contact. Performance in these settings must describe the composition of the selected set as well as separation between outcome classes.

AUROC describes separation through the probability that a randomly drawn positive outranks a randomly drawn negative. It remains unchanged when prevalence changes and the class-conditional score distributions remain fixed. Selected-set purity behaves differently. Suppose a threshold selects $80\%$ of positive cases and $5\%$ of negative cases. In a balanced population of 10,000 cases, the selected set contains 4,000 positives and 250 negatives, giving precision $0.94$. In a population with $1\%$ positives, it contains 80 positives and 495 negatives, giving precision $0.14$. The threshold has the same true-positive and false-positive rates in both populations. The number of negative cases available for selection changes the meaning of its precision.

Average precision aggregates precision as the threshold is lowered. It therefore inherits the prevalence of the population being evaluated. Two studies can evaluate identical class-conditional ranking behavior and report different AP values because their populations contain different proportions of positive cases. Comparisons across studies require a common population interpretation.

Some applications have one well-defined target prevalence. Others serve several populations, settings, or planning scenarios whose prevalences differ. Selecting one convenient prevalence hides the populations in which a performance claim may weaken. Averaging over prevalences is the other obvious move, and it requires a distribution over target populations. Where such a distribution has a defensible scientific basis it should be used, and prevalence then belongs with the other sampled quantities discussed below. Where it does not, an assumed distribution lets strong performance in one population compensate for weak performance in another. A declared set of target prevalences makes the more demanding claim. The reported level must remain defensible throughout that set.

Replication supplies a second source of challenge. A fixed model can be tested on new evaluation units. A claim intended for new environments must also face environmental variation. A claim about a development procedure must survive new training, tuning, and model-selection runs. An evaluation-generating regime declares which of these mechanisms may vary. Independent executions of that regime give a performance claim repeated opportunities to fail.

These are two parts of one thing. A performance claim is stated against a declared \emph{challenge space}, the set of conditions under which it must survive. The target-prevalence set $\PSet$ declares the population conditions in that space. The regime $\Regime$ declares the data-generating and development conditions. The replication count $K$ declares how many independent executions must retain the claim. Together they give the resulting statement its scope, and a value quoted without them states nothing in particular.


The available predictions determine how the statement enters a decision. One model evaluated on held-out outcomes can be located relative to the chance level of its finite evaluation design. Compatible reports from separate studies then support separate design-relative interpretation, and not an ordering. Competing models evaluated on the same units permit a paired analysis of the advantage retained throughout $\PSet$. The paired quantity supplies performance evidence for a choice or replacement decision.

\section{A common average-precision effect scale}
\label{sec:common-scale}

This section develops the first of the three declarations, the target-prevalence set $\PSet$. Sections~\ref{sec:regime} and \ref{sec:support} develop the regime $\Regime$ and the replication count $K$.

\subsection{Prior-standardized average precision}

At score threshold $c$, let
\[
r(c)=\Prb(S\geq c\mid Y=1),
\qquad
f(c)=\Prb(S\geq c\mid Y=0).
\]
These are the true-positive and false-positive rates. Precision in a population with prevalence $\pi\in(0,1)$ is
\begin{equation}
p_\pi(c)
=
\frac{\pi r(c)}
{\pi r(c)+(1-\pi)f(c)}.
\label{eq:reference-precision}
\end{equation}

Precision odds show how the population prior and ranking behavior enter:
\begin{equation}
\frac{p_\pi(c)}{1-p_\pi(c)}
=
\frac{\pi}{1-\pi}
\frac{r(c)}{f(c)}.
\label{eq:precision-odds}
\end{equation}
The first factor gives the target-population prior odds. The second is a threshold-specific likelihood ratio determined by the class-conditional score distributions. Prior standardization changes the first factor and preserves the second. The same factorization supports prior-probability and cost adjustment in classification \cite{elkan2001foundations}. It is also used to evaluate precision-based metrics at a fixed reference class ratio \cite{siblini2020calibration,degeest2016online}.

Applying Equation~\eqref{eq:reference-precision} across the distinct score thresholds gives
\[
\AP_\pi(D),
\]
the prior-standardized AP of evaluation $D$ at target prevalence $\pi$. The empirical definition and tie convention appear in Appendix~\ref{app:ap-definition}.

\subsection{The populations the transformation reaches}
\label{sec:transport}

Equation~\eqref{eq:reference-precision} holds $r(c)$ and $f(c)$ fixed and moves $\pi$. It therefore describes target populations that differ from the observed one in the proportion of positive cases and in nothing else. This is the assumption the whole construction rests on, and it can be stated in a form the reader can check.

Let $X$ be the case features, $Y$ the outcome, and $S$ the model score. The transformation requires one invariance, and it can be stated without any graph:
\begin{equation}
P(S\geq c\mid Y=1)
\ \text{ and }\
P(S\geq c\mid Y=0)
\ \text{ agree in the observed and target populations.}
\label{eq:selection-diagram}
\end{equation}
The rates $r(c)$ and $f(c)$ measured in the observed population are then the rates in the target, and Equation~\eqref{eq:reference-precision} evaluates the target exactly. Only the class prior differs. This is the score-level form of prior-probability shift, or label shift. The literature ordinarily assumes the stronger invariance of $P(X\mid Y)$, which implies Equation~\eqref{eq:selection-diagram} for any fixed scoring rule but is not implied by it \cite{saerens2002adjusting,storkey2009training,lipton2018detecting}.

It is convenient to picture the condition as a diagram in which one node marks what differs between the two populations and sends a single arrow into $Y$, none into $X$ or $S$ \cite{pearl2011transportability,bareinboim2013transportability}. The picture is an aid rather than the warrant. Reading the joint law as $P(Y)\,P(X\mid Y)$ is a factorization and not a claim about the mechanism by which cases arise, and in many predictive settings the features precede the outcome. Equation~\eqref{eq:selection-diagram} is what the analysis uses, and it is a statement about two distributions.

Two departures may defeat the transformation. If the class-conditional feature distribution differs, cases carrying the same label do not look the same in the two populations, which is what covariate shift and concept shift both produce. This moves the rates only where the features that changed reach the score; a shift confined to features the fitted model does not use leaves $P(S\mid Y)$ where it was, and Equation~\eqref{eq:selection-diagram} can hold in spite of it. If the measurement or scoring process differs, the rates move with no change in the cases at all. Where either departure does move the rates, Equation~\eqref{eq:reference-precision} reweights rates that do not describe the target population, and no choice of $\pi$ recovers it.

Equation~\eqref{eq:selection-diagram} concerns which mechanisms are held constant rather than anything about the observed data, so it is settled before the analysis and belongs with the prespecified items of Section~\ref{sec:reporting}. Section~\ref{sec:separability} gives the consequence when it fails.

\subsection{Chance normalization}

Population random ranking selects the same fraction of each class, so $r(c)=f(c)$ and $p_\pi(c)=\pi$. Its prior-standardized AP is therefore $\pi$. Define
\begin{equation}
\CNAP_\pi(D)
=
\frac{\AP_\pi(D)-\pi}{1-\pi}.
\label{eq:cnap}
\end{equation}
CNAP is zero for population random ranking, one for perfect ranking, and negative below chance. Its range at prevalence $\pi$ is
\[
\left[-\frac{\pi}{1-\pi},1\right].
\]

The normalization has the general form
\[
\frac{\text{observed}-\text{chance}}
{\text{ceiling}-\text{chance}},
\]
which is familiar from chance-corrected agreement statistics \cite{cohen1960coefficient}. Earlier work has rescaled AP using either the observed prevalence or a fixed class ratio \cite{suyuan2015relationship,bestgen2021fisher}. Equation~\eqref{eq:cnap} first evaluates AP in the declared target population and then uses that population's random-ranking value as the lower anchor. Every target prevalence consequently shares the endpoints zero and one.

At empirical threshold $h$,
\begin{equation}
\frac{p_{\pi,h}-\pi}{1-\pi}
=
\frac{\pi(r_h-f_h)}
{\pi r_h+(1-\pi)f_h}.
\label{eq:threshold-cnap}
\end{equation}
Smaller prevalences give false positives greater influence because negative cases are more abundant. Larger prevalences credit class separation over a broader recall range. Model ordering can change with $\pi$ even after chance normalization.

\subsection{The target-prevalence challenge set}

Let
\[
\PSet\subset(0,1)
\]
be a nonempty compact set of target prevalences chosen from scientific knowledge, deployment plans, or a benchmark specification. A closed interval $[\pi_L,\pi_U]$ is a natural choice when every prevalence between two plausible limits is relevant. A finite set is appropriate when the application concerns named operating populations.

\begin{definition}[Prevalence-robust chance-normalized average precision]
\begin{equation}
\RobCNAP(D)
=
\inf_{\pi\in\PSet}\CNAP_\pi(D).
\label{eq:robust-cnap}
\end{equation}
\end{definition}

The quantity is the largest level retained at every prevalence in the declared set. The empirical CNAP profile is continuous on the compact set $\PSet$, so its infimum is attained. A minimizing prevalence may be recorded as
\[
\pi^\dagger(D)
\in
\arg\min_{\pi\in\PSet}\CNAP_\pi(D).
\]
Several minimizing values may exist.

When $\PSet=\{\pi_\star\}$, Equation~\eqref{eq:robust-cnap} reduces to the single-reference-prevalence effect. If $\PSet$ has maximum $\pi_U<1$, a common range is
\[
\RobCNAP(D)
\in
\left[-\frac{\pi_U}{1-\pi_U},1\right].
\]

The pointwise profile
\[
\pi\longmapsto\CNAP_\pi(D)
\]
shows which target populations limit the result and where model ordering changes. The profile serves as a diagnostic. Equation~\eqref{eq:robust-cnap} supplies the scalar claim.

The set must be declared independently of the observed performance profile. Its width controls the range of population conditions allowed to challenge the claim. If $\Pi_1\subseteq\Pi_2$, then
\[
\underline{\CNAP}_{\Pi_2}(D)
\leq
\underline{\CNAP}_{\Pi_1}(D).
\]
Widening the set cannot improve the retained level.

\subsection{Where the worst case occurs}
\label{sec:worst-case}

Appendix~\ref{app:ap-definition} shows that a threshold contribution is nondecreasing in $\pi$ whenever $r_h\geq f_h$. A model whose contributing thresholds all lie above the ROC diagonal therefore has a nondecreasing CNAP profile, and
\[
\inf_{\pi\in[\pi_L,\pi_U]}\CNAP_\pi(D)
=
\CNAP_{\pi_L}(D).
\]
Prevalence-robust performance then reduces to performance at the lower endpoint, and the declared set enters only through that endpoint.

Two consequences follow. For a single model, the robust scalar may carry little beyond a well-chosen point prevalence, and the choice of $\pi_L$ then requires the scientific justification that would otherwise attach to a single reference prevalence. For a paired comparison, the difference of two nondecreasing profiles need not be monotone. The limiting prevalence can fall in the interior and can differ across replications. The construction contributes most where model ordering changes across the target set.

How often interior minima occur is an empirical question. Recording the observed minimizer in each report allows it to be assessed across applications. Section~\ref{sec:paired} returns to this point, because the paired comparison is where a nonmonotone difference makes the declared set do work that no single prevalence could do.

\section{Performance across replications}
\label{sec:regime}

The second declaration is the regime. A performance claim acquires meaning from the conditions under which it could fail. These conditions are represented by an evaluation-generating regime.

One replication is produced by a fixed set of mechanisms acting in order. An environment $E$ is drawn. Training data are drawn within it and a development procedure returns a fitted model $\widehat f$. Evaluation units $U$ are drawn, a measurement process gives them features $X$, outcomes $Y$ are realized, and $\widehat f$ assigns scores:
\begin{equation}
\begin{gathered}
E\longrightarrow D_{\mathrm{train}}\longrightarrow\widehat f,
\qquad
E\longrightarrow U\longrightarrow(X,Y),
\qquad
(\widehat f,X)\longrightarrow S,
\\
D^{\mathrm{rep}}=\{(S_i,Y_i)\}_{i\in U}.
\end{gathered}
\label{eq:regime-graph}
\end{equation}

\begin{definition}[Evaluation-generating regime]
An evaluation-generating regime $\Regime$ is a probability law for one complete replication of the claim under study. It is specified by declaring, for each node of Equation~\eqref{eq:regime-graph}, whether that node is redrawn from its mechanism or held at the value the observed evaluation gave it.
\end{definition}

Declaring a regime is therefore a choice of where to cut Equation~\eqref{eq:regime-graph}. Holding a node freezes everything upstream of it and leaves everything downstream free to vary. Three common cuts illustrate the range of claims.

\begin{table}[ht]
\centering
\caption{Illustrative evaluation-generating regimes. Each entry records the status of one node of Equation~\eqref{eq:regime-graph}.}
\label{tab:regimes}
\begin{tabularx}{\textwidth}{@{}>{\raggedright\arraybackslash}p{0.25\textwidth}XXX@{}}
\toprule
Node & Fixed-model & New-environment & Repeated-development \\
\midrule
Evaluation units $U$ & redrawn & redrawn & redrawn \\
Environment $E$ & held & redrawn & redrawn \\
Measurement process $X\mid U$ & held & declared & declared \\
Replication class counts & declared & declared & declared \\
Fitted model $\widehat f$ & held & held & redrawn \\
\bottomrule
\end{tabularx}
\end{table}

Every held node is a scientific conjecture, and the graph says which one. Holding the measurement process asserts that the mechanism generating $X$ from $U$ behaves comparably in the populations the claim covers. Holding $\widehat f$ confines the claim to one realized model. Redrawing $\widehat f$ moves the claim from a model to the procedure that produced it, and costs a complete rerun of that procedure in every replication. External replication can expose failures in these conjectures, which is what makes them conjectures rather than definitions.

Write $n_+^{\mathrm{obs}}$ and $n_-^{\mathrm{obs}}$ for the observed evaluation counts. Let $n_+^{\mathrm{rep}}$ and $n_-^{\mathrm{rep}}$ denote the counts generated in one replication. A same-design analysis equates the replication counts with the observed counts. A prospective analysis declares the counts planned for future evaluations. The observed counts govern the information available for estimation. The replication counts help define the claim.

\subsection{Two parts of one challenge space}
\label{sec:challenge-space}

The set $\PSet$ is held fixed across replications. Every generated evaluation is examined at every prevalence in that set, prevalence being varied through Equation~\eqref{eq:reference-precision}. It is not drawn as an additional random variable.

Prevalence is not a different kind of challenge from the ones the regime carries. Both declare conditions under which the claim must hold. The difference is that prevalence is the only such condition that can be evaluated from a realized evaluation without generating a new one, and Section~\ref{sec:transport} says what buys that. The invariance of Equation~\eqref{eq:selection-diagram} leaves the score mechanism untouched, so the threshold rates $(r_h,f_h)$ of a single evaluation are the rates of every target population in $\PSet$, and Equation~\eqref{eq:precision-odds} converts them into $\CNAP_\pi$ at every $\pi$ at once. What a model would have done in a different environment, under a different measurement process, or after a different development run is not fixed by any invariance the observed evaluation supplies. It has to be drawn.

The challenge space therefore has a computed part and a simulated part:
\[
\underbrace{\PSet}_{\text{evaluated by infimum}}
\qquad
\underbrace{\Regime,\ K}_{\text{explored by sampling}}.
\]
Worst case is taken over the whole space; expectation is taken over the sampled part alone. This is what Equation~\eqref{eq:robust-support} expresses, and it is why the infimum sits inside the replication rather than outside it.

The two remain separate in the notation, and $S_{K,\PSet}^{\mathrm{AP}}(\Regime)$ keeps them in distinct argument positions, because they are combined by different operators. Folding $\PSet$ into $\Regime$ would suggest averaging over prevalence, which is the construction Section~\ref{sec:motivation} declined.

\subsection{When separability fails}
\label{sec:separability}

Analytic separability is not free. It is Equation~\eqref{eq:selection-diagram} under another description, and the boundary between the computed and simulated parts of the challenge space falls exactly where that invariance fails.

Where the invariance fails, prevalence does not leave the challenge space. It migrates from the computed part to the simulated part. The rates $r(c)$ and $f(c)$ then vary with $\pi$, no reweighting recovers the target, and the environment component of the regime must generate the accompanying change directly. The declared claim is unaffected. Only the machinery that interrogates it moves, and the cost is that the analysis now requires data from the populations in question.

This gives a test for the boundary of the method, and the test is a question about mechanisms rather than about the observed profile. Ask whether the populations named by $\PSet$ are reached from the observed one by changing the class prior alone. A benchmark whose evaluation set was built by subsampling negatives from a single pool qualifies, since the class-conditional distributions are identical across prevalences by construction. A specialist referral population set against a primary-care population does not, since the two differ in case severity as well as in prevalence, and severity moves the class-conditional feature distribution. Where Equation~\eqref{eq:selection-diagram} holds, the prevalence challenge is computed and costs nothing. Where it does not, it must be simulated, and Appendix~\ref{app:regime} gives the construction.

Let
\[
D^{\mathrm{rep}}\sim P_{\Regime}
\]
denote one execution of the regime and define
\begin{equation}
Z_{\PSet}
=
\RobCNAP(D^{\mathrm{rep}}).
\label{eq:robust-replication-effect}
\end{equation}
The law of $Z_{\PSet}$ is the prevalence-robust replication distribution. Its center describes expected worst-prevalence performance. Its shape describes how that retained level varies across replications.

Independence applies to complete replications. Observations within a replication may be dependent. Clustered, matched, longitudinal, and multi-environment designs require resampling at the coarsest unit the design treats as independent, the cluster or matched set or site rather than the individual observation. That unit is called the \emph{resampling unit} throughout. A repeated-development regime reruns the complete fitting pipeline. These designs are admitted throughout the single-model development below. The paired criterion of Section~\ref{sec:paired} imposes a separate condition, on the arrangement of outcome labels within one evaluation rather than on the unit of resampling, and Section~\ref{sec:paired-scope} states it.

Replication-based assessment of classifier performance has precedent in work on generalization error and the replicability of learning experiments \cite{nadeau2003inference,bouckaert2004estimating}. The present replication distribution concerns a prevalence-robust CNAP effect generated under a declared regime.

\section{Replication-survival support}
\label{sec:support}

The third declaration is the replication count. This section builds the operator that consumes it, then combines all three declarations into the primary quantity.

Let $T$ be a scalar performance effect whose larger values indicate better performance. An evaluation yielding $t$ retains every claim of the form $T\geq s$ with $s\leq t$. Values above $t$ are defeated by that evaluation.

Suppose two independent replications yield $0.70$ and $0.05$. A claim based on both results cannot exceed $0.05$. The larger value supplies no compensation for a level defeated by the other replication. The greatest jointly retained level is their minimum.

For $K$ independent executions of $\Regime$, let
\[
T_1,\ldots,T_K
\iid
P_{\mathrm{rep}}^{\Regime,T}.
\]
Define
\begin{equation}
S_K(T;\Regime)
=
\E_{\Regime}
\left[
\min_{1\leq j\leq K}T_j
\right].
\label{eq:support-operator}
\end{equation}
The operator averages the greatest level retained by each set of $K$ replications. It remains on the scale of $T$.

This construction expresses corroboration through survival under independent opportunities for defeat \cite{popper1959logic}. The regime determines the kinds of challenges presented to the claim. The value of $K$ determines how many independent challenges it faces. The default $K=2$ is the smallest value that requires joint survival.

The word \emph{supported} is used here in its Popperian sense. A level is supported when it has been exposed to declared opportunities for refutation and has not been defeated by them. Corroboration in this sense is not a probability, is not a degree of belief, and does not accumulate into confirmation; Popper was explicit that a well-corroborated claim has earned no probability of being true, only a record of what it has withstood. Read this way, $K$ specifies how many independent executions of $\Regime$ the claim is required to survive. It is declared by the analyst rather than counted in the data, and the estimator of Section~\ref{sec:estimation} projects that criterion from the evidence one evaluation supplies; its computational resamples are not themselves empirical replications. How demanding the required executions are is settled by the regime and not by $K$. A fixed-model regime redraws evaluation units and holds everything else, so its replications expose the claim to sampling variation alone, and a large $K$ under a narrow regime is a long sequence of easy tests.

Severity in the error-statistical tradition is instead computed from the sampling distribution of a test statistic and reports how probable a less favorable result would have been had the claim been false \cite{mayo1996error}. Equation~\eqref{eq:support-operator} returns no such probability. The one quantity here with an error-statistical character is the permutation $p$-value of Equation~\eqref{eq:permutation-value}, reported separately and for a narrower purpose.

The operator follows from the criterion rather than being selected among candidates, and the sense in which it does can be stated exactly. Four requirements are natural for any summary of a replication distribution meant to be read on the scale of the effect: it depends on the distribution alone, respects first-order stochastic dominance, returns $c$ for a claim that retains $c$ in every replication, and is additive across components that rise and fall together. These place the summary in the distortion family of Appendix~\ref{app:families}, where it weights each level by an increasing function $g$ of the probability that one replication retains it \cite{schmeidler1986integral,yaari1987dual}. Joint survival then fixes $g$. A level that one replication retains with probability $u$ is retained by all $K$ with probability $u^K$, so $g(u)=u^K$; and the powers are the only distortions under which the weight on a jointly retained level is the product of the weights on the separately retained ones, which is what rules the alternatives out. Appendix~\ref{app:characterization} gives the argument.

That fifth requirement is what separates the operator from its neighbors. A lower quantile and an expected shortfall satisfy the first four and fail it, since their distortions are not multiplicative. They are not weaker answers to the question of joint survival. They answer different questions, in which no requirement of joint survival appears. Appendix~\ref{app:proofs} gives a frequency-qualified retained level for readers who want one of them answered.

If $T$ has lower bound $L$ and upper bound $U$, then
\begin{equation}
S_K(T;\Regime)
=
L+\int_L^U
\Prb_{\Regime}(T>s)^K\,ds.
\label{eq:support-survival}
\end{equation}
At each level $s$, the powered survival probability is the probability that all $K$ replications retain the claim. Appendix~\ref{app:proofs} gives the derivation.

At $K=2$, the minimum identity gives
\begin{equation}
S_2(T;\Regime)
=
\E_{\Regime}[T]
-
\frac12
\E_{\Regime}
\left[
|T_1-T_2|
\right].
\label{eq:k2-decomposition}
\end{equation}
Supported performance equals expected performance across replications minus half their expected absolute disagreement. The second term records the reduction in the retained level produced by variation across the regime.

\subsection{Prevalence-robust supported performance}

Apply Equation~\eqref{eq:support-operator} to the within-replication effect in Equation~\eqref{eq:robust-replication-effect}.

\begin{definition}[Prevalence-robust supported CNAP]
\begin{equation}
\SRob
=
\E_{\Regime}
\left[
\min_{1\leq j\leq K}
\inf_{\pi\in\PSet}
\CNAP_\pi(D_j)
\right].
\label{eq:robust-support}
\end{equation}
\end{definition}

The supported level survives every target prevalence in $\PSet$ and every one of the $K$ replications. The order of operations carries this interpretation. Each replication first receives the lowest value attained over the prevalence set. Replication support is then applied to those retained values.

Collapsing the two operations shows what the definition does to the challenge space of Section~\ref{sec:challenge-space}:
\[
\SRob
=
\E_{\Regime}
\left[
\inf_{\substack{\pi\in\PSet\\1\leq j\leq K}}
\CNAP_\pi(D_j)
\right].
\]
This is the split of Section~\ref{sec:challenge-space} applied to the primary quantity, and Section~\ref{sec:infimum-inside} shows what happens when the two operations are taken in the wrong order.
If $L_{\PSet}$ is a common lower bound, its joint-survival form is
\begin{equation}
\SRob
=
L_{\PSet}
+
\int_{L_{\PSet}}^1
\Prb_{\Regime}
\left\{
\CNAP_\pi(D)>s
\text{ for every }\pi\in\PSet
\right\}^{K}
ds.
\label{eq:robust-joint-survival}
\end{equation}
The integrand gives the probability that every declared population condition retains level $s$ in all $K$ replications. Writing the complement makes the reading direct. Let
\begin{equation}
\RefSet_s
=
\left\{
(\pi,j)\in\PSet\times\{1,\ldots,K\}
:
\CNAP_\pi(D_j)\leq s
\right\}
\label{eq:refutation-class}
\end{equation}
collect the declared conditions that defeat level $s$. The level survives exactly when $\RefSet_s$ is empty, and Equation~\eqref{eq:robust-joint-survival} becomes the following, with the independence across replications now carried inside the event rather than shown as a power:
\begin{equation}
\SRob
=
L_{\PSet}
+
\int_{L_{\PSet}}^{1}
\Prb_{\Regime}
\left(
\RefSet_s=\emptyset
\right)
ds .
\label{eq:refutation-integral}
\end{equation}
Supported performance is the integral over levels of the probability that nothing in the declared challenge space defeats that level. The prevalence set and the replication count enter through one object, the index set $\PSet\times\{1,\ldots,K\}$ that $\RefSet_s$ is drawn from. Section~\ref{sec:severity} takes them together on that basis.

At $K=2$,
\begin{equation}
S_{2,\PSet}^{\mathrm{AP}}(\Regime)
=
\E_{\Regime}[Z_{\PSet}]
-
\frac12
\E_{\Regime}
\left[
|Z_{\PSet,1}-Z_{\PSet,2}|
\right].
\label{eq:robust-k2}
\end{equation}
The first term is expected worst-prevalence performance. The second records disagreement between the worst-prevalence levels attained in independent replications. Their minimizing prevalences may differ.

\subsection{Why the infimum belongs inside each replication}
\label{sec:infimum-inside}

Faced with a prevalence set and a replication criterion, the natural construction is to apply them in the order they were introduced. Calculate supported performance separately at each prevalence, giving the pointwise supported profile
\[
\pi
\longmapsto
S_K(\CNAP_\pi;\Regime),
\]
then take its lowest value,
\[
\inf_{\pi\in\PSet}S_K(\CNAP_\pi;\Regime).
\]
This is what a reader of a prevalence profile computes without thinking about it, and it is what a profile of separately supported values reports.

It is the weaker quantity. Every replication is challenged at the same prevalence before the prevalence minimum is taken, so a replication that fails badly at $\pi=0.01$ and a replication that fails badly at $\pi=0.40$ are never both charged for their own worst case. Only one prevalence is ever held against the whole array. Equation~\eqref{eq:robust-support} instead lets each realized replication reveal the target population that defeats it. The ordering is
\begin{equation}
S_K\left(
\inf_{\pi\in\PSet}\CNAP_\pi;
\Regime
\right)
\leq
\inf_{\pi\in\PSet}
S_K(\CNAP_\pi;\Regime).
\label{eq:strong-weak-order}
\end{equation}
The two sides agree when one prevalence minimizes the profile almost surely. They separate as the limiting prevalence moves between replications. Appendix~\ref{app:proofs} gives the short proof and shows that the entire gap comes from taking the expectation after the prevalence search rather than before.

The same reversal appears before any replication support is applied:
\begin{equation}
\E_{\Regime}
\left[
\inf_{\pi\in\PSet}\CNAP_\pi
\right]
\leq
\inf_{\pi\in\PSet}
\E_{\Regime}
\left[
\CNAP_\pi
\right].
\label{eq:prevalence-selection}
\end{equation}
Two named costs follow, and the paper uses these names throughout. The \emph{prevalence-challenge cost}
\[
C_{\mathrm{prev}}
=
\inf_{\pi\in\PSet}\E_{\Regime}[\CNAP_\pi]
-
\E_{\Regime}[Z_{\PSet}]
\]
is the difference between the two sides of Equation~\eqref{eq:prevalence-selection}. It is what the claim gives up by allowing the limiting prevalence to differ from one replication to the next, and it is zero when a single prevalence minimizes the realized profile almost surely. The \emph{replication-disagreement cost}
\[
C_{\mathrm{rep}}
=
\E_{\Regime}\left[Z_{\PSet}\right]
-
\SRob
\]
is what the claim then gives up by requiring $K$ replications to retain the same level. Together they give the reading used in Section~\ref{sec:reporting}:
\begin{equation}
\SRob
=
\underbrace{\inf_{\pi\in\PSet}\E_{\Regime}[\CNAP_\pi]}_{\text{anchor}}
-
C_{\mathrm{prev}}
-
C_{\mathrm{rep}}.
\label{eq:two-costs}
\end{equation}
The anchor is expected performance at the least favorable prevalence for the average profile. Both costs are nonnegative, the second because $\min_{j\leq K}Z_{\PSet,j}\leq Z_{\PSet,1}$ in every realization, and both are available from the bootstrap array at no extra cost. The replication cost rises with $K$, since requiring more replications to agree gives up more.

At $K=2$ the minimum identity evaluates it in closed form,
\[
C_{\mathrm{rep}}
=
\frac12
\E_{\Regime}
\left[
|Z_{\PSet,1}-Z_{\PSet,2}|
\right],
\]
half the expected disagreement between two replications, which is the second term of Equation~\eqref{eq:robust-k2}. No expression in two moments is available for $K>2$. Appendix~\ref{app:penalty} gives the derivation and the general form.

\subsection{Severity under prevalence and replication challenges}
\label{sec:severity}

A test is made more severe by giving the claim more ways to fail, which is the sense in which a bolder claim is a better one \cite{popper1963conjectures}. Equation~\eqref{eq:refutation-integral} says what that amounts to here. Enlarge the index set from which $\RefSet_s$ is drawn, and the probability that $\RefSet_s$ is empty cannot rise. Two of the three declarations enlarge it by inclusion, and they do so in the same way.

\begin{proposition}[Challenge monotonicity]
\label{prop:challenge-monotonicity}
Let $\Pi_1\subseteq\Pi_2$ and $K_1\leq K_2$. Then
\begin{equation}
S_{K_2,\Pi_2}^{\mathrm{AP}}(\Regime)
\leq
S_{K_1,\Pi_1}^{\mathrm{AP}}(\Regime).
\label{eq:challenge-monotonicity}
\end{equation}
\end{proposition}

The proof is one line and appears in Appendix~\ref{app:proofs}. Widening $\PSet$ and raising $K$ both enlarge $\PSet\times\{1,\ldots,K\}$, so both can only add elements to $\RefSet_s$, at every level and in every realization. The prevalence set introduces more population conditions that can defeat the claim and the replication count introduces more independent occasions on which it can be defeated. These are one operation performed on the two factors of a single index set, which is why they need one proof and not two.

The regime is the third declaration and it behaves differently. It does not enlarge the index set. It changes the law of the values that set indexes. Regimes are not ordered by inclusion, so Proposition~\ref{prop:challenge-monotonicity} has no counterpart for them. What is available is an ordering by dispersion: a regime that spreads the replication distribution while holding its mean fixed lowers the supported value, by the concave-order property in Appendix~\ref{app:families}. Declaring that measurement conditions or development runs may vary is a severity choice of the same kind, resting on a different argument and holding only between regimes that the dispersion ordering compares.

Severity in each direction is a claim about the raw supported value, for the reason Section~\ref{sec:calibration-def} gives.


The replication-survival operator applies to any scalar effect oriented so that larger values indicate better performance. A loss $L(D)$ can be represented by $T(D)=-L(D)$. A paired loss comparison can use
\[
T(D)=L_B(D)-L_A(D),
\]
so positive values favor model $A$.

The prevalence challenge generalizes on the criterion of Section~\ref{sec:challenge-space} rather than on prevalence itself. Any dimension of a challenge space that can be evaluated from one realized evaluation, without generating another, can be handled by infimum at no sampling cost. Decision costs and case-mix weights are candidates wherever an analogue of Equation~\eqref{eq:precision-odds} isolates them in a single factor. Dimensions without that structure have to be simulated, and they belong in the regime.

\section{Estimation from one observed evaluation}
\label{sec:estimation}

The regime-indexed quantity in Equation~\eqref{eq:robust-support} is a population functional. One observed evaluation supplies a data-conditional representation of its replication distribution. The representation must reflect the mechanisms permitted to vary by the declared regime.

\subsection{Canonical fixed-model estimator}

Suppose $D_{\mathrm{obs}}$ is an independent or held-out evaluation of a fixed fitted model. Separate its positive and negative evaluation units into
\[
D_{\mathrm{obs},+},
\qquad
D_{\mathrm{obs},-},
\]
with counts $n_+^{\mathrm{obs}}>0$ and $n_-^{\mathrm{obs}}>0$. The canonical same-design estimator sets the replication counts equal to these observed counts.

For each bootstrap replication $b=1,\ldots,B$:
\begin{enumerate}
\item Sample $n_+^{\mathrm{rep}}$ units with replacement from $D_{\mathrm{obs},+}$.
\item Sample $n_-^{\mathrm{rep}}$ units with replacement from $D_{\mathrm{obs},-}$.
\item Combine the sampled units into $D^{(b)}$.
\item Evaluate the complete CNAP profile over $\PSet$.
\item Record
\begin{equation}
z_b
=
\inf_{\pi\in\PSet}
\CNAP_\pi(D^{(b)}).
\label{eq:bootstrap-robust-effect}
\end{equation}
\end{enumerate}

The array $z_1,\ldots,z_B$ is the empirical prevalence-robust replication distribution. Its construction holds the fitted model fixed and redraws evaluation units within outcome class. Hierarchical designs replace unit-level resampling with the resampling unit declared by the regime.

For $K=2$, the all-pairs estimator is
\begin{equation}
\widehat S_{2,\PSet,B}^{\,\mathrm{AP}}
=
\frac{2}{B(B-1)}
\sum_{1\leq i<j\leq B}
\min(z_i,z_j).
\label{eq:all-pairs-estimator}
\end{equation}
It is the order-two U-statistic with kernel $\min(z_i,z_j)$ \cite{hoeffding1948ustat}. It uses every unordered pair of independent computational replications. The equivalent sorted-array calculation appears in Appendix~\ref{app:computation}. The subscript $B$ records the number of computational replications and is written explicitly wherever the Monte Carlo layer is at issue; where it is not, the same estimator is written $\SRobHat$.

For general $K\leq B$, the corresponding U-statistic is
\begin{equation}
\widehat S_{K,\PSet,B}^{\,\mathrm{AP}}
=
\binom{B}{K}^{-1}
\sum_{\substack{I\subset\{1,\ldots,B\}\\|I|=K}}
\min_{i\in I}z_i.
\label{eq:general-u-statistic}
\end{equation}
An order-statistic formula avoids explicit enumeration.

The observed prevalence-robust effect is
\[
z_{\mathrm{obs}}
=
\inf_{\pi\in\PSet}
\CNAP_\pi(D_{\mathrm{obs}}).
\]
The observed minimizer $\pi^\dagger(D_{\mathrm{obs}})$ helps explain which target population limits the point estimate. The minimizing prevalence may vary across bootstrap replications, so it does not replace the full within-replication optimization.

\subsection{Continuous intervals and finite grids}

When $\PSet=[\pi_L,\pi_U]$, the empirical CNAP profile is continuous on a compact interval and attains its minimum. A one-dimensional optimizer can evaluate Equation~\eqref{eq:bootstrap-robust-effect}. The optimization tolerance should be small relative to the reported precision.

A prespecified finite grid
\[
\Pi_{\star,G}=\{\pi_1,\ldots,\pi_G\}
\]
defines a finite challenge set. Its minimum is exact for that set. If the grid approximates a continuous interval, report the grid and assess the approximation error. A coarse grid can miss an interior minimum and raise the reported robust value.

Every prevalence within a replication shares the same resampled units. This preserves the joint profile and allows the limiting prevalence to emerge from that replication. Independent resampling at separate prevalences would destroy this dependence.

\subsection{Prospective and repeated-development regimes}

A prospective fixed-model analysis replaces the observed bootstrap sizes with declared future evaluation counts and draws from the empirical class-conditional score pools. This changes the replication target and should be labeled as a prospective design.

Claims about new environments require data from multiple environments and a resampling hierarchy that redraws them. Claims about development procedures require complete reruns of fitting, tuning, and model selection. A bootstrap of one fixed score vector identifies neither target.

\section{Conditional calibration for one model}
\label{sec:calibration}

This section takes up the first of the two tasks, interpreting one model considered on its own. Section~\ref{sec:paired} takes up the second.

CNAP assigns zero to population random ranking. A finite evaluation design does not. Three mechanisms move its chance level away from zero, and they do not move it the same way. Stepwise AP has upward finite-sample bias under random ranking \cite{bestgen2015random}, which raises it. The replication-disagreement cost $C_{\mathrm{rep}}$ lowers it. The prevalence-challenge cost $C_{\mathrm{prev}}$ lowers it further. Only the complete pipeline determines where the three leave it, which is why the null is computed by running that pipeline rather than by correcting the value afterwards.

\subsection{A design with no signal at all}

The smallest case makes the point without any estimation. Take an evaluation with two observations, one positive and one negative, and let $\PSet=\{1/2\}$. There is no model worth speaking of and no association to find. Under random ranking the positive observation ranks first half the time and second half the time. Stepwise AP is then $1$ or $1/2$, so CNAP is $1$ or $0$.

The population random-ranking value is zero, but the replication mean is
\[
\E[Z_{\PSet}]
=
\frac12,
\]
and for two independent replications
\[
S_{2,\PSet}^{\mathrm{AP}}(\Regime_0)
=
\E[\min(Z_{\PSet,1},Z_{\PSet,2})]
=
\frac14.
\]
A procedure applied to pure noise returns a supported value of $0.25$ on a scale whose zero was supposed to mean chance. Nothing has gone wrong. The discreteness of stepwise AP puts half the mass at CNAP $=1$, and no amount of chance normalization can remove a distortion that lives in the finite design rather than in the population.

Two observations is an extreme case, and the effect shrinks as $n_+$ grows. It does not vanish, it interacts with ties and with $K$, and widening $\PSet$ adds the prevalence search to the same calculation. There is no general formula for where the three mechanisms leave the chance level of a particular design. The remedy is to run the whole procedure on outcomes that carry no signal and see what it returns.

The size of the displacement is a property of the score vector, not only of the sample size. Equation~\eqref{eq:empirical-ap} evaluates precision at the thresholds where positive cases enter, so it never samples a threshold with no positive case above it, and the precision values it averages are drawn from a favorable subset of those available. This mechanism requires distinct thresholds to be available. A score vector with a single tie block offers none: its only threshold has $r_h=f_h=1$, prior-standardized precision equals $\pi$ exactly, and $\AP_\pi=\pi$ with no variability at all. A vector of distinct scores offers the maximum. The finite-design chance level therefore rises with the granularity of the score vector, and two models evaluated on the very same units can sit at different chance levels for reasons unconnected to their merit. Section~\ref{sec:paired-anchor} shows what this does to a comparison between them.

\subsection{Conditional fixed-prediction null}

Suppose a developed model has been evaluated on outcomes that played no role in its development. Let $\mathbf{s}$ be its realized score vector and let $\mathbf{y}$ contain the observed positive and negative counts.

\begin{definition}[Conditional fixed-prediction null]
Condition on the score vector, its tie structure, the evaluation units, and the class counts. Assign outcomes to scores through the permutations allowed by the evaluation design.
\end{definition}

For independent evaluation units, every assignment preserving the class counts is allowed. Clustered and matched designs use restricted permutations at their resampling unit. The restriction is harmless here, because the null is conditional on the realized score vector and locates one model against the chance level of its own design. Section~\ref{sec:paired-scope} explains why the paired result of Section~\ref{sec:paired} cannot be extended to restricted designs in the same way. The null asks how much prevalence-robust supported performance the procedure produces after the association between the realized scores and outcomes has been removed. Permutation tests for fixed classifier predictions provide the corresponding conditional no-association test \cite{ojala2010permutation}.

For each allowed or sampled permutation $m=1,\ldots,M$:
\begin{enumerate}
\item Attach the permuted outcomes to the fixed scores.
\item Construct $B$ bootstrap replications under the declared fixed-model regime.
\item Within every bootstrap replication, calculate
\[
z_b^{(m)}
=
\inf_{\pi\in\PSet}
\CNAP_\pi(D^{(m,b)}).
\]
\item Apply the support estimator to $z_1^{(m)},\ldots,z_B^{(m)}$ and record
\[
\widehat S_{K,\PSet,B}^{\,\mathrm{AP},(m)}.
\]
\end{enumerate}

The permutation replicates support two summaries, and the pair is easy to confuse.

The first is their mean, the estimated robust conditional chance level,
\begin{equation}
\SRobNullHat
=
\frac{1}{M}
\sum_{m=1}^{M}
\widehat S_{K,\PSet,B}^{\,\mathrm{AP},(m)}.
\label{eq:robust-null}
\end{equation}
This is the anchor that Section~\ref{sec:calibration-def} calibrates against.

The second is their upper tail,
\begin{equation}
p
=
\frac{
1+
\#\left\{
m:
\widehat S_{K,\PSet,B}^{\,\mathrm{AP},(m)}
\geq
\SRobHat
\right\}
}{M+1}.
\label{eq:permutation-value}
\end{equation}
This is a permutation $p$-value for the conditional no-association null. It is conditional on the realized score vector and its tie structure, so it refers to this evaluation rather than to any population, and it carries the usual cautions attaching to any $p$-value.

The two summaries use different features of the same $M$ replicates. Equation~\eqref{eq:robust-null} uses only their mean, so the calibrated value of Section~\ref{sec:calibration-def} discards the spread of the null distribution entirely. Equation~\eqref{eq:permutation-value} is the only quantity in the paper that uses that spread. Two designs with the same null mean and very different null dispersion produce identical calibrated values and quite different $p$, which is why both belong in a report.

The prevalence infimum, bootstrap support calculation, and permutation average form one procedure. Taking the infimum of separately calibrated pointwise values would use a different null construction and a different scale.

\subsection{Conditionally calibrated robust support}
\label{sec:calibration-def}

Perfect separation produces $\CNAP_\pi=1$ at every prevalence. Under the canonical stratified bootstrap, every bootstrap replication therefore has $z_b=1$ and the supported value is one.

\begin{definition}[Conditionally calibrated prevalence-robust supported CNAP]
\begin{equation}
\SRobCalHat
=
\frac{
\SRobHat-\SRobNullHat
}{
1-\SRobNullHat
}.
\label{eq:robust-calibration}
\end{equation}
\end{definition}

The calibrated score assigns the robust conditional chance level zero and perfect separation one. A value at or below the null anchor supplies no positive calibrated claim. Report $\SRobHat$ and $\SRobNullHat$ alongside the calibrated result, since a ratio cannot be checked without its numerator and denominator, and report the permutation $p$-value of Equation~\eqref{eq:permutation-value} for the dispersion the calibrated value omits.

Set inclusion has a direct severity interpretation for the raw quantity in Equation~\eqref{eq:robust-support}. Calibration changes both the observed supported value and its null anchor. The calibrated score can therefore move in either direction when $\PSet$ is widened. Statements about increased prevalence severity should refer to the raw supported value.

Two models can have different tie structures and different conditional null anchors, so their calibrated scales differ and two calibrated values can be neither subtracted nor ranked against each other. Nor is the raw paired difference free of the problem: differing anchors displace it from zero by an amount neither model's calibration records. Section~\ref{sec:paired} removes that term by adopting a tie convention under which the two anchors coincide, having first shown that estimating a correction for it fails.

The conditional null assumes that the outcomes used for calibration were isolated from model development. Reuse of evaluation outcomes for tuning, feature selection, stopping, or model choice invalidates the fixed-prediction permutation. A repeated-development claim requires a null procedure for the complete development process.

\section{Paired model comparison and replacement decisions}
\label{sec:paired}

The second task is a decision between two models evaluated together. The quantity it reports is
\begin{equation}
\theta_{A:B}
=
S_K(\RobDta;\Regime)
=
\E_{\Regime}
\left[
\min_{1\leq j\leq K}
\inf_{\pi\in\PSet}
\left\{
\CNAPta_{A,\pi}(D_j)
-
\CNAPta_{B,\pi}(D_j)
\right\}
\right],
\label{eq:robust-paired-support}
\end{equation}
the advantage of $A$ over $B$ that survives every prevalence in $\PSet$ and every one of $K$ replications of $\Regime$. Apart from the tie convention marked by the tilde, which Section~\ref{sec:paired-convention} introduces, this is the construction of Sections~\ref{sec:regime} and~\ref{sec:support} applied to a difference rather than to a level. It is the declared estimand, and what follows establishes what can and cannot be said about it.

Four things have to be settled. A paired report cannot be assembled from two separate reports, which Section~\ref{sec:marginal} shows. It must carry no positive anchor contributed by the evaluation design alone, which Sections~\ref{sec:paired-anchor} to \ref{sec:reference-law} take up and Section~\ref{sec:paired-scope} bounds. It must not be able to favor both models at once, which is Section~\ref{sec:orientation}. And it must carry from the observed evaluation to the population the decision concerns, which Section~\ref{sec:replacement} states.

\subsection{Why separate reports do not compose}
\label{sec:marginal}

The obvious economy is to evaluate each model on its own and subtract. Write
\[
A_{\PSet}(D)
=
\inf_{\pi\in\PSet}\CNAPta_{A,\pi}(D),
\qquad
B_{\PSet}(D)
=
\inf_{\pi\in\PSet}\CNAPta_{B,\pi}(D)
\]
for the marginal robust values. For every evaluation $A_{\PSet}(D)-B_{\PSet}(D)\geq\RobDta(D)$, and at every $K$
\begin{equation}
S_K(A_{\PSet};\Regime)
-
S_K(B_{\PSet};\Regime)
\geq
S_K(A_{\PSet}-B_{\PSet};\Regime)
\geq
S_K(\RobDta;\Regime).
\label{eq:robust-marginal-order}
\end{equation}
Appendix~\ref{app:proofs} proves both for any pair of prevalence-indexed profiles, so the chain holds under either tie convention. The first step charges for each marginal value meeting its own worst replication, the second for each model meeting the prevalence worst for it alone. What model $A$ reports is the level it retains at $A$'s worst prevalence in $A$'s worst replication, and what $B$ reports is the same for $B$; their difference compares two worst cases that never occurred together. Equation~\eqref{eq:robust-paired-support} instead asks whether $A$ leads at the same prevalence in the same replication, and takes the worst of that.

Differencing two separately reported values therefore overstates the paired advantage, and the chain bounds the error in one direction only. It is an algebraic property of the infimum and the minimum, assuming nothing about the design and using no property of average precision, so no convention repairs it and no design escapes it. Separate studies also cannot recover the joint replication behavior Equation~\eqref{eq:robust-paired-support} requires. The rest of this section takes up a second obstacle, which does depend on how average precision treats equal scores.

\subsection{Pairing does not remove the design anchor}
\label{sec:paired-anchor}

Suppose models $A$ and $B$ are evaluated on the same units. At prevalence $\pi$ the natural comparison is the paired effect
\begin{equation}
\Delta_\pi(D)
=
\CNAP_{A,\pi}(D)
-
\CNAP_{B,\pi}(D),
\label{eq:paired-effect}
\end{equation}
with positive values favoring model $A$, and the natural summary is its infimum over $\PSet$. Both models meet the same units and the same outcomes, so the difficulty of the realized evaluation is common to them and cancels from the difference. This subsection shows that something else does not cancel, and that Equation~\eqref{eq:paired-effect} is therefore not the quantity to report.

Section~\ref{sec:calibration} established that a finite evaluation design displaces the chance level of CNAP, and that the displacement depends on the granularity of the score vector. It is a property of a model's own scores rather than of the shared data, so two models compared on the same units contribute two displacements to their difference. Writing $\beta_A$ and $\beta_B$ for the expected values the two models attain under a design with no association between scores and outcomes,
\[
\E[\Delta_\pi]
=
\beta_A-\beta_B
\]
in that design. Pairing removes what the two models share and leaves this difference behind. It is guaranteed to be zero when the two score vectors have the same tie structure, and need not be zero otherwise.

The smallest case shows the size of the effect. Take four evaluation units, two of them positive, and let $\PSet=\{1/2\}$. Model $A$ assigns four distinct scores whose ordering has nothing to do with the outcomes. Model $B$ assigns every unit the same score. Neither model carries any signal.

Under the six equally likely assignments of two positive labels to four units, Equation~\eqref{eq:empirical-ap} gives $\AP_{1/2}(A)$ equal to $1$, $5/6$, $3/4$, $7/12$, $1/2$ and $5/12$, so $\E[\AP_{1/2}(A)]=49/72$. Model $B$ has a single threshold block with $r=f=1$, so $\AP_{1/2}(B)=1/2$ in every assignment and its CNAP is exactly zero. Converting through Equation~\eqref{eq:cnap},
\[
\E\left[\Delta_{1/2}\right]
=
\frac{13}{36}
\approx
0.361 .
\]
The replication operator does not repair this. Let $\Regime_0$ denote a law under which the outcome assignment is redrawn and nothing else varies, so that a replication is one of the six assignments. Then
\[
S_2(\RobDelta;\Regime_0)
=
\frac{29}{216}
\approx
0.134 ,
\]
and the procedure reports a supported advantage for a model that differs from its competitor only in breaking ties arbitrarily.

Two features of this matter for what follows. The sign of the distortion follows the finer score vector, so the direction of the error depends on which model happens to be more granular. And the replication-disagreement cost of Section~\ref{sec:infimum-inside} sometimes offsets the distortion and sometimes compounds it, so the null value of the raw paired quantity under the block convention cannot be signed in advance. Appendix~\ref{app:paired-null} gives these computations together with larger cases in which neither model is degenerate and cases in which the coarser model is favored. The situations concerned are ordinary ones, such as a model with discrete or heavily rounded outputs compared against one producing a continuous score.

Two routes lead out. The displacement can be measured under a joint permutation and subtracted, which is what Section~\ref{sec:calibration} does for an isolated model, or it can be removed by definition. The first route fails, and Appendix~\ref{app:paired-null} works the case: the displacement is not an additive constant, since it must vanish where both models approach perfect separation and there is no room for it, so a correction calibrated where the models are uninformative over-corrects where they are not. Carried through, such a correction can report each of two identically perfect models as superior to the other. The second route is taken next.

\subsection{Tie-neutral average precision}
\label{sec:paired-convention}

The displacement is not a fact about average precision. It is a fact about the convention adopted in Appendix~\ref{app:ap-definition} for handling equal scores, under which a tie block contributes the value obtained by treating it as a single threshold. A different convention removes the displacement outright, and it is the convention this paper requires for paired comparison.

\begin{definition}[Tie-averaged average precision]
Let $\APta_\pi(D)$ denote the average of Equation~\eqref{eq:empirical-ap} over the rankings consistent with the observed score vector, each tie block contributing the mean of the values obtained under its internal orderings. Write $\CNAPta_\pi$ for the corresponding chance normalization through Equation~\eqref{eq:cnap}, and
\[
\Dta_\pi(D)
=
\CNAPta_{A,\pi}(D)-\CNAPta_{B,\pi}(D),
\qquad
\RobDta(D)
=
\inf_{\pi\in\PSet}\Dta_\pi(D).
\]
\end{definition}

The convention has a direct reading. Let $U$ be a vector of independent continuous draws and let $\varepsilon$ be small enough that adding $\varepsilon U$ to the scores reorders units only within tie blocks. Then
\begin{equation}
\APta_\pi(s,y)
=
\E_U\left[\AP_\pi(s+\varepsilon U,y)\mid s,y\right].
\label{eq:jitter}
\end{equation}
Tie-averaged AP is the expected value of ordinary AP once the ties have been broken at random. Tied cases are ordered arbitrarily because the model has supplied no basis for ordering them, and averaging removes the dependence on which arbitrary ordering was drawn rather than committing to one of them. Appendix~\ref{app:tie-computation} evaluates the expectation in closed form, so no randomization is performed and the convention costs no enumeration.

\subsection{Neutrality to uninformative refinement}
\label{sec:refinement}

What the averaging removes is credit for breaking ties in a particular way. What it leaves untouched is credit for ordering units correctly. The following states it exactly, and it is what answers Section~\ref{sec:paired-anchor}.

Let $\Rank_A$ be the strict ranking obtained from $A$'s scores with any ties $A$ itself retains resolved uniformly at random, which is the ranking $\APta_{A,\pi}$ averages over. Stating the hypothesis on $\Rank_A$ covers the case in which $A$ refines only some of $B$'s blocks or leaves ties of its own.

\begin{proposition}[Refinement neutrality]
\label{prop:refinement}
Suppose $A$ agrees with $B$ on the ordering of $B$'s tie blocks, and that conditionally on $S_B$ and $Y$ the ranking $\Rank_A$ is uniform over the strict rankings consistent with $S_B$. Then for every $\pi$,
\[
\E\left[\APta_{A,\pi}\mid S_B,Y\right]
=
\APta_{B,\pi}.
\]
\end{proposition}

The proof is two lines and appears in Appendix~\ref{app:tie-convention}: averaging $\AP_\pi$ over a uniform draw from the rankings consistent with $S_B$ is what defines $\APta_{B,\pi}$.

Three things follow. The statement conditions on the realized outcomes, so it holds when the models carry signal and not only when they do not: granularity unrelated to the outcomes earns nothing in expectation, however good the ranking it was added to. It asks nothing of the design either, so it survives the matched, clustered and heterogeneous-risk evaluations that Section~\ref{sec:paired-scope} excludes from the result below. And it is an expectation result and nothing stronger: one arbitrary tie-breaking can align with the labels by chance and show a realized advantage, and what the convention removes is the positive expected credit for having broken the ties at all.

A model whose finer scores genuinely separate the units within a block still earns the separation, and one that orders them wrongly is still charged for it. The portion of a comparison that arbitrary tie-breaking would have generated has zero conditional expectation, and it has zero conditional expectation at every target prevalence.

\subsection{Reference behavior under exchangeable labels}
\label{sec:reference-law}

Proposition~\ref{prop:refinement} concerns two models related by refinement. A second result covers two arbitrary models carrying no association with the outcomes, and it needs a reference law to state.

\begin{definition}[Exchangeable-label reference law]
\label{def:label-condition}
A reference law $P_0$ for one evaluation assigns
\[
P_0\left(Y=y\mid S_A,S_B,n_+\right)
=
\binom{n}{n_+}^{-1}
\]
for every label vector $y$ containing $n_+$ positives.
\end{definition}

Three things have to be kept apart. The regime $\Regime$ is the law the claim is about, and under it the scores may carry as much association with the outcomes as they can. The reference law $P_0$ removes that association by randomizing the labels conditionally on both score vectors. The design requirement, stated in Section~\ref{sec:paired-scope}, is what makes $P_0$ a defensible description of a signal-free version of this study rather than an arbitrary construction. No claim is made that the observed evaluation is drawn from $P_0$, and an informative model guarantees that it is not. The word exchangeability is used in a different sense in the causal literature, where it names the comparability of treated and untreated groups.

\begin{proposition}[Uniformity of the design anchor]
\label{prop:uniformity}
Fix $n_+$, $n_-$ and $\pi$. Under $P_0$, $\E_0[\APta_\pi]$ takes the same value for every score vector, whatever its tie structure, and consequently $\E_0[\Dta_\pi]=0$.
\end{proposition}

The proof is two lines and appears in Appendix~\ref{app:tie-convention}. For a fixed ranking of the units, $P_0$ makes the law of the label sequence read along that ranking the same whatever the ranking, so the expected value of Equation~\eqref{eq:empirical-ap} depends on $n_+$, $n_-$ and $\pi$ alone and not on which ranking was fixed. Tie-averaged AP is a convex combination of stepwise values over the rankings consistent with the score vector, and a convex combination of a constant is that constant.

Write $\Regime_0$ for the law that follows $\Regime$ in drawing environments, evaluation units, measurement conditions and fitted models, with the replication class counts unchanged, and assigns the outcome labels of each evaluation according to Definition~\ref{def:label-condition} conditionally on everything drawn before them.

\begin{proposition}[Reference behavior of the paired criterion]
\label{prop:conservative}
Under $\Regime_0$, $\E_{\Regime_0}[\Dta_\pi]=0$ for every $\pi\in\PSet$, and consequently
\[
S_K(\RobDta;\Regime_0)
\leq
\E_{\Regime_0}[\RobDta]
\leq
0 ,
\]
with the same bound for the reversed comparison $B$ against $A$.
\end{proposition}

The conditioning is what admits the regimes of Section~\ref{sec:regime}. Writing $E$ for everything drawn before the labels of one evaluation,
\[
\E_{\Regime_0}\left[\Dta_\pi\right]
=
\E\left\{
\E_0\left[\Dta_\pi\mid E\right]
\right\}
=
0
\]
by Proposition~\ref{prop:uniformity} applied within each $E$. A regime that redraws environments, measurement conditions or the fitted model is therefore not excluded, and the infimum over $\PSet$ and the minimum over replications then carry the bound downward.

What this establishes should be stated narrowly. Two models with different tie structures cannot produce a positive supported advantage under the reference law: the difference between them is not anchored away from zero by the granularity of their scores. It does not establish that $\theta_{A:B}$ is centered at zero whenever two informative models are equally good. There is no general statistical null corresponding to equally good rankings, since two models may differ in where over $\PSet$ they perform well, in their replication variability, and in their dependence on one another while agreeing in marginal average precision. Section~\ref{sec:replacement} accordingly builds the decision on the estimand and an outer bound rather than on a test of equality, and no paired $p$-value accompanies it. Permuting the outcomes would remove the association for both score vectors at once, describing a pair of uninformative models rather than a pair of equally informative ones, so two models that were genuinely good and genuinely equal would sit far in its upper tail.

Both inequalities run downward in both orientations, so a design meeting the requirement of Section~\ref{sec:paired-scope} does not induce a positive reference anchor, and can conceal a real advantage rather than supply one. The bound is on the population functional: a finite estimate can still exceed zero under $\Regime_0$ through sampling variability, which is what the outer layer of Section~\ref{sec:replacement} addresses. That is the price of removing the design term by definition rather than estimating it, and Section~\ref{sec:paired-anchor} gave the reason for paying it: a quantity that can be removed by definition should not be estimated, and an estimated correction would carry its own error into the decision.

For estimation, the same resampled units and outcomes are used for both models within bootstrap replication $b$. Calculate
\[
\delta_b
=
\inf_{\pi\in\PSet}
\left\{
\CNAPta_{A,\pi}(D^{(b)})
-
\CNAPta_{B,\pi}(D^{(b)})
\right\},
\]
then apply the all-pairs estimator of Equation~\eqref{eq:all-pairs-estimator} to $\delta_1,\ldots,\delta_B$. A minimizing value $\pi_\Delta^\dagger(D)\in\arg\min_{\pi\in\PSet}\Dta_\pi(D)$ identifies a target population that limits the comparison. At $K=2$,
\begin{equation}
\theta_{A:B}
=
\E_{\Regime}[\RobDta]
-
\frac12
\E_{\Regime}
\left[
|\underline{\widetilde{\Delta}}_{\PSet,1}-\underline{\widetilde{\Delta}}_{\PSet,2}|
\right].
\label{eq:robust-paired-k2}
\end{equation}

\subsection{Scope and honest evaluation requirements}
\label{sec:paired-scope}

Two conditions bound Proposition~\ref{prop:conservative}. The first concerns the arrangement of labels within an evaluation, the second what the evaluation outcomes were allowed to influence. Proposition~\ref{prop:refinement} needs neither.

\paragraph{Designs the reference law does not describe.}
Section~\ref{sec:regime} admits clustered, matched, longitudinal and multi-environment designs, and Section~\ref{sec:calibration} assigns outcomes through the permutations such designs permit. Definition~\ref{def:label-condition} asks for more than those permutations supply. Restricted randomization removes the association between scores and outcomes without making every arrangement of the labels equally likely. A ranking can therefore be null with respect to the outcomes and still be aligned with the strata the design randomizes within, and tie averaging does not touch that alignment. What Equation~\eqref{eq:tie-averaged-ap} removes is credit for breaking ties inside a block. What it cannot remove is credit for arranging the blocks in agreement with the randomization strata.

Take four units in two matched pairs, each pair containing exactly one positive, so that the permitted permutations are the independent label swaps within pairs, and let $\PSet=\{1/2\}$. Let model $A$ give every unit the same score and let model $B$ place the first pair above the second, so that the tie blocks of $B$ are the strata themselves. Every block of $B$ then contains exactly one positive under every permitted assignment, and $A$ has a single block, so $\Dta_{1/2}$ takes the same value under all four assignments:
\[
\Dta_{1/2}\equiv\frac{1}{36},
\qquad
C_{\mathrm{rep}}(\RobDta)=0,
\qquad
S_K(\RobDta;\Regime_0)=\frac{1}{36}
\quad\text{for every }K .
\]
Neither model carries any association with the outcomes, and one is reported as supportedly superior to the other. Raising $K$ does not damp the advantage, because there is no replication disagreement to charge for, and an outer bound around a degenerate quantity does not remove it either. The value decays with evaluation size, slowly, and is still $0.008$ at $n=200$; it does not decay with severity, and severity is the mechanism this paper relies on. Appendix~\ref{app:tie-convention} gives the scaling, together with a second counterexample that is tie-free in both models and defeats Proposition~\ref{prop:uniformity} without defeating Proposition~\ref{prop:conservative}.

\paragraph{The resulting restriction.}
Proposition~\ref{prop:conservative} is available where Definition~\ref{def:label-condition} is a defensible reference law for the design, which asks that the conditional label distribution be invariant under the full symmetric group once the signal is removed. This holds for i.i.d. held-out units. It is a statement about the signal-free law and not about a randomization anyone performed; outside a randomized experiment the design supplies the ground for asserting the invariance rather than the invariance itself. Independence alone does not supply it: units that are independent but not identically distributed, as when their baseline risks differ, need not make every arrangement equally likely, and Proposition~\ref{prop:uniformity} asks for exactly that. Sampling design and protocol carry most of the justification, together with the judgment that the signal-free outcome law is exchangeable across units. Within one replication that judgment is defensible when no dependence joins the outcomes of distinct evaluation units, either directly or through a parent they share that the replication does not redraw, and no unit-level parent of $Y$ varies across units in a way the score vector tracks. The two failures already described are one of each: matched pairs share a parent within a pair, which is what lets the tie blocks of model $B$ coincide with the strata, and differing baseline risk is a unit-level parent of $Y$ that a score vector can track. Where the law is not defensible, $\theta_{A:B}$ remains computable and remains a description of what one model retained over another; what is unavailable is the assurance that the evaluation design contributes no positive paired anchor to it. The tie-free counterexample of Appendix~\ref{app:tie-convention} shows that the failure is not confined to differing granularity, since a ranking aligned with the randomization strata produces an anchor with no tie structure involved at all. Appendix~\ref{app:tie-convention} describes the obstruction exactly, through the group of label permutations the design randomizes over, and a convention averaging over an orbit of that group would be the natural repair. It is left open here.

\paragraph{Honest evaluation.}
Proposition~\ref{prop:conservative} describes $\theta_{A:B}$ under $\Regime_0$. It says nothing about whether an evaluation carried out under some other law is informative about $\theta_{A:B}$ at all. Both score vectors, the orientation of the comparison, the prevalence set $\PSet$, the replication count $K$, the tie convention and the margin $d$ must be fixed without using the evaluation outcomes. This excludes using those outcomes for training or tuning, for early stopping or feature selection, for selecting the two models from a larger set, for deciding which model is designated $A$, and for choosing $\PSet$, $K$ or $d$.

It is tempting to think the paired comparison escapes this, since its reference law is free of contamination by construction. It does not. Suppose model $A$ used the evaluation labels and model $B$ did not. The observed statistic favors $A$, every bootstrap replication of it favors $A$, and a reference law in which the leakage is absent removes nothing from either. Proposition~\ref{prop:conservative} constrains the behavior of the statistic under $\Regime_0$; it does not make a contaminated $\Regime$ informative about an uncontaminated advantage. Where evaluation outcomes reached development, the estimand is the performance of the leaky procedure, and it is worth reporting only if that is what the decision concerns.

A repeated-development regime is admitted and is not an exception. Each development run uses the training data drawn for it, which is the content of the claim being made about the procedure; what it may not use is the evaluation outcomes against which the run is then assessed, and those must have been unavailable to both procedures. Cross-validation and cross-fitting sit outside the statement, since withholding a unit's own label does not prevent the remaining evaluation labels from reaching its prediction or from selecting the pipeline that produced it. Within these bounds the comparison reaches a setting the isolated report does not. The permutation of Section~\ref{sec:calibration} conditions on a realized score vector that a repeated-development regime redraws with the model, so Section~\ref{sec:reporting} withdraws the calibrated value there, while Definition~\ref{def:label-condition} randomizes labels under the regime rather than around a fixed vector and Proposition~\ref{prop:conservative} continues to hold.

\paragraph{What survives the restriction.}
The coherence bound of Section~\ref{sec:orientation} is an algebraic property of the minimum and holds under any design. Proposition~\ref{prop:refinement} conditions on the realized outcomes and asks nothing of the design either, so the neutrality to arbitrary refinement that motivated the convention is not lost on a matched or clustered evaluation. The non-composition result of Section~\ref{sec:marginal} holds under any design. And the argument of Section~\ref{sec:paired-anchor} against a centered correction does not use the reference law, so a restricted design does not reopen that route.

\subsection{Coherence and the no-verdict region}
\label{sec:orientation}

A procedure that compares two models must not be able to report that each of them beats the other. This one cannot, and the reason needs nothing beyond the shape of the replication operator, which is not odd. Since $\min_j x_j\leq\max_j x_j$ pointwise and $S_K(-T;\Regime)=-\E_{\Regime}[\max_{j\leq K}T_j]$,
\begin{equation}
S_K(\RobDta;\Regime)
+
S_K(-\RobDta;\Regime)
=
-\E_{\Regime}
\left[
\max_{j\leq K}\underline{\widetilde{\Delta}}_{\PSet,j}-\min_{j\leq K}\underline{\widetilde{\Delta}}_{\PSet,j}
\right]
\leq
0 .
\label{eq:orientation}
\end{equation}
Both orientations cannot show a supported advantage at once, since that would require the expected range of the replications to be negative. The bound follows from the minimum alone, so unlike Propositions~\ref{prop:uniformity} and~\ref{prop:conservative} it requires no reference law and holds under any evaluation design, including the designs Section~\ref{sec:paired-scope} has just excluded.

The same algebra says the comparison can decline to choose. Both orientations may be at or below zero, and the interval between them is the region in which the replication criterion returns no verdict. At $K=2$ its width is
\[
-\left\{
S_2(\RobDta;\Regime)+S_2(-\RobDta;\Regime)
\right\}
=
\E_{\Regime}
\left|
\underline{\widetilde{\Delta}}_{\PSet,1}
-
\underline{\widetilde{\Delta}}_{\PSet,2}
\right|
=
2C_{\mathrm{rep}}(\RobDta),
\]
exactly twice the replication-disagreement cost of Section~\ref{sec:infimum-inside}.

Section~\ref{sec:severity} gives the reading. A no-verdict result reports that the declared challenges defeated the claim of superiority in each direction, which is what a test unable to separate two models should report. The width is the resolving power of the comparison, and it narrows when the evidence for a difference stabilizes across the regime and in no other way. Both orientations should be computed, and a report that gives only one has answered half the question it was asked.

\subsection{Replacement criterion}
\label{sec:replacement}

Let $d\geq0$ be the smallest retained CNAP advantage that would justify replacing model $B$ with model $A$. The performance criterion is
\begin{equation}
\theta_{A:B}>d.
\label{eq:replacement}
\end{equation}
The choice $d=0$ requires a positive advantage throughout the prevalence set under the replication criterion. A positive $d$ states a practical threshold. Where the same CNAP difference carries different practical meaning at different prevalences, a declared margin function $d(\pi)$ can replace the constant, and Appendix~\ref{app:variable-margin} gives that form.

The criterion is stated on the raw tie-averaged scale with no anchor subtracted, so $d$ is not silently inflated or discounted by a design term. What $d$ means is a retained advantage on the tie-neutral comparison scale, and that is narrower than a deployment quantity. By Equation~\eqref{eq:jitter} the scale scores a tie block as though its members had been ordered at random. A deployment that selects whole blocks, or treats a discrete risk category as operationally indivisible, does something else, and where that is so the blockwise performance of both models should be reported alongside the comparison. The paired statistic answers whether one ranking is better once credit for arbitrary tie order has been removed; whether the improvement is worth having at the operating point in use is a separate question, settled by thresholds and costs. What the analyst gives up is stated in Section~\ref{sec:orientation}: the criterion is conservative, and a real advantage smaller than the replication disagreement will not be declared.

The fourth requirement is what the outer layer supplies. Proposition~\ref{prop:conservative} locates the reference behavior of $\theta_{A:B}$ under $\Regime_0$ and says nothing about its value under $\Regime$, nor about any one estimate of either, so a finite bootstrap estimate can exceed zero through estimation variability alone and a positive point estimate does not by itself establish superiority. The proposed inferential criterion for replacement is that an outer lower bootstrap bound lie above $d$. Appendix~\ref{app:calibration} states what remains to be established about that bound, which is offered as a procedure rather than as an interval endpoint of known coverage. It is a further challenge rather than a different kind of warrant: none of the three declarations exposes the claim to the possibility that the observed evaluation is an unrepresentative draw from its regime, and a claim that clears $d$ at the point estimate and fails at the bound has been defeated by it. Costs, feasibility, fairness, probability calibration, and operational consequences can also enter the final decision.


\section{Interpretation, uncertainty, and reporting}
\label{sec:reporting}

\subsection{What the robust value carries}

The prevalence-robust supported value remains on the CNAP scale. It records a level retained across the population conditions in $\PSet$ and the replication challenges in $\Regime$. Its value depends on the replication counts, tie structure, class-conditional score distributions, and sources of variation represented by the regime.

Evaluation size enters through the replication distribution, and it enters through both costs at once. Section~\ref{sec:sparse} works through the case of few positive cases, where the effect is largest.

The robust value does not carry a fixed confidence level or recurrence probability. It does not protect against population changes outside $\PSet$. It also does not protect against changes in class-conditional score behavior unless the regime represents them.

One limitation is not a matter of what the regime represents. Every quantity here describes a ranking in a population generating outcomes on its own, and the applications of Section~\ref{sec:motivation} do not leave the population alone. Selecting a case for review or for contact is an action taken on that case, and the action can change the outcome that would otherwise have been observed \cite{perdomo2020performative}. No regime built from Equation~\eqref{eq:regime-graph} reaches this, since a regime describes populations that differ in how they were drawn rather than populations that have been acted upon. A supported claim is a claim about ranking behavior, and it becomes a claim about deployed performance only where the selection leaves outcomes undisturbed.

\subsection{Sparse positive cases}
\label{sec:sparse}

Stepwise AP moves in increments of $1/n_+$. When positive cases are few, one rank change moves the profile at every prevalence, and three effects follow.

The replication distribution of $Z_{\PSet}$ becomes coarse with a heavy lower tail, so $C_{\mathrm{rep}}$ grows. The minimizing prevalence becomes unstable across replications, so $C_{\mathrm{prev}}$ grows as well. The conditional null rises, because the finite-sample upward bias of stepwise AP under random ranking is largest when $n_+$ is small.

Raw robust support therefore falls, while the calibrated value can move in either direction because its anchor also rises. Reporting $C_{\mathrm{prev}}$ and $C_{\mathrm{rep}}$ alongside the bootstrap distribution of $\pi^\dagger$ shows which mechanism is responsible. Neither a larger bootstrap count nor a narrower prevalence set recovers information the evaluation does not contain.

\subsection{Monte Carlo and estimation uncertainty}

The inner bootstrap approximates the data-conditional support functional. Its Monte Carlo error can be estimated from the U-statistic influence values in Appendix~\ref{app:computation}. The permutation average has a second Monte Carlo component governed by $M$.

An outer resampling layer addresses how well the observed evaluation identifies the supported target. Each outer sample must repeat the complete procedure:
\[
\text{outer resample}
\longrightarrow
\text{inner replication array}
\longrightarrow
\text{prevalence infima}
\longrightarrow
\text{support functional}.
\]
For calibrated performance, the robust conditional null must also be recomputed inside every outer sample. For paired superiority no null is recomputed, since Section~\ref{sec:paired-convention} estimates no correction; both models simply remain paired throughout the outer and inner layers.

The infimum can make the estimator nonsmooth when the minimizing prevalence changes. Sparse outcomes and flat minima can further affect interval behavior. Coverage should be studied by simulation under designs resembling the intended application.

\subsection{Reporting}

A report has three parts: what survived, what it was declared to survive, and what explains the number. The first is a sentence. The second is prespecified. The third belongs in the supplement. Producing every quantity in this paper for every analysis would obscure the claim it was built to make.

\subsubsection*{What survived}

For an isolated model the reported quantity is the calibrated value $\SRobCalHat$. The raw supported value and the conditional null travel with it and are never separated from it, because a ratio cannot be checked without its numerator and denominator. For a decision between two models the reported quantity is $S_K(\RobDta;\Regime)$ on the raw tie-averaged CNAP scale, together with the value obtained under the reversed orientation. A single sentence carries either claim:

\begin{quote}
Across the prevalence range $[\pi_L,\pi_U]$ and $K$ independent replications of the declared regime, the model achieved a calibrated prevalence-robust CNAP of [value] (raw supported level [value] against a finite-design chance level of [value]). The worst case occurred at prevalence [value].
\end{quote}

The permutation $p$-value follows in the next sentence rather than inside this one. The quoted claim is a statement about magnitude, and a $p$-value placed within it invites the whole claim to be read as a significance result.

\subsubsection*{What it was declared to survive}

The prespecified declarations are $\PSet$ with the scientific basis for its boundaries, the AP convention and the invariance of Equation~\eqref{eq:selection-diagram} that licenses prevalence transport, the regime with its class counts and resampling unit, $K$, the replacement margin $d$ where one applies, the separation between model development and evaluation outcomes, and for a paired comparison the reference law of Definition~\ref{def:label-condition}, stated for the units of one evaluation and justified against the design criterion of Section~\ref{sec:paired-scope}. The last of these is not the resampling unit under another name; a design may declare a matched set as its resampling unit and still leave that reference law indefensible. These are design, not results. They are written once, before the analysis, and they belong in the methods rather than mixed among the estimates.

\subsubsection*{What explains the number}

Table~\ref{tab:reporting} places each remaining quantity against the question it answers.

\begin{table}[ht]
\centering
\caption{Which quantity answers which question.}
\label{tab:reporting}
\begin{tabularx}{\textwidth}{@{}XXl@{}}
\toprule
Question & Quantity & Where \\
\midrule
How much survived? & $\SRobHat$ & headline \\
How far is that above this design's chance level? & $\SRobCalHat$ & headline \\
Could the null pipeline alone have produced this? & permutation $p$-value, Equation~\eqref{eq:permutation-value} & next sentence \\
Does $A$ beat $B$? & $S_K(\RobDta;\Regime)$, both orientations & headline \\
Does the claim clear a threshold? & proposed outer lower bound & only if a threshold exists \\
Why is the value low? & anchor, $C_{\mathrm{prev}}$, $C_{\mathrm{rep}}$, spread of $\pi^\dagger$ & supplement \\
Has the comparison reached a verdict? & the no-verdict width, Equation~\eqref{eq:orientation} & with the headline \\
Which populations limit it? & the CNAP profile & supplement \\
Did the computation converge? & Monte Carlo errors & methods \\
\bottomrule
\end{tabularx}
\end{table}

The outer layer is the item most often reported out of habit. It is required when a claim must clear a threshold, whether a replacement margin $d$ or an absolute minimum acceptable level, since $S_K(\RobDta;\Regime)>d$ and $\SRob>c$ take the same treatment. Absent a threshold, a two-sided interval reports estimation uncertainty and a one-sided bound answers a question nobody asked. The distinction has a practical edge: an outer resample costs one inner bootstrap for paired superiority and a factor of $M$ more for the calibrated value, since the conditional null must be recomputed inside every outer sample. The comparison estimates no correction, and so pays for none.

The profile uses the same prespecified set and shared replications as the scalar calculation.

\subsection{How to read the two headline values}

The raw supported value states how much performance survived. Use it to compare severity settings and to test a margin. The calibrated value states how far that level exceeds the mean chance level of this evaluation design. Use it to interpret one model. Both are magnitudes on a performance axis. Neither is a probability, and neither is a confidence bound.

The permutation $p$-value is a probability, and it answers the narrower question the calibrated value cannot. Calibration shifts and rescales by the null mean alone, so a supported value can sit above that mean and still fall well within the range the null pipeline produces. It is conditional on the realized score vector, so it speaks about this evaluation and not about the population. It is a subordinate diagnostic and not the claim.

Two cautions attach to the calibrated value and to neither of the others. A wider $\PSet$ moves its anchor as well as its numerator, so a claim that a wider set was more severely tested must cite the raw value. And two calibrated values are neither differenced nor ranked against each other: each subtracts its own model's null anchor and divides by its own model's remaining range, and two such affine maps reverse an ordering as readily as they distort a difference, so a model with the higher raw supported value and the higher anchor can sit below its rival once both are calibrated. A calibrated value locates one model against the chance level of its own design, and that is the whole of what it does. The paired quantity escapes both cautions. It is reported on the raw tie-averaged scale with no anchor subtracted, so it falls monotonically as $\PSet$ widens, and Section~\ref{sec:paired-convention} removes the need for an anchor rather than assembling one.

The calibrated value is unavailable where the fixed-prediction permutation is invalid, which includes repeated-development regimes and any case where evaluation outcomes informed model development. The headline then reverts to the raw supported value, labeled as uncalibrated. A paired comparison remains available in a repeated-development regime, where the permutation is the wrong construction rather than an invalidated one. It does not remain available where evaluation outcomes reached model development, and Section~\ref{sec:paired-scope} states what an honest paired evaluation requires of both models.


\subsection{What an empirical assessment must establish}

The framework is developed here without an empirical illustration, and that is its principal remaining limitation. The questions below must be settled before the robust scalar replaces a single-prevalence report in practice.

Four of them test the paired construction directly. With no signal in either model, varying the tie granularity should show the block convention manufacturing an advantage of the size Section~\ref{sec:paired-anchor} predicts, across evaluation sizes, prevalences and degrees of asymmetry, while the tie-averaged supported advantage stays nonpositive on average. Beginning instead from an informative coarse model and splitting its tie blocks at random should show no expected tie-averaged advantage at any prevalence, which is the empirical content of Proposition~\ref{prop:refinement} and the case that separates it from Proposition~\ref{prop:uniformity}. Introducing label-related ordering inside those same blocks should show the comparison crediting information genuinely added. And the matched and heterogeneous-risk designs of Section~\ref{sec:paired-scope} should be reproduced across a range of sizes, to establish what the restriction costs, how slowly the artifact decays, and whether any subclass can be readmitted under a design-specific convention.

The operating characteristics follow from these. Power and no-verdict frequency are one quantity by Section~\ref{sec:orientation}, and what has to be established is how large a genuine advantage must be, relative to the replication disagreement, before the criterion declares it. Outer intervals require a coverage study when the minimizer sits at an endpoint, in the interior, at several points, and when it moves under small perturbations. Monte Carlo stability requires the number of replications at which the supported value settles under prevalence optimization.

The remaining questions concern the single-model quantity and whether the declared prevalence set changes any scientific conclusion. Matched evaluations with similar observed profiles should be compared to see whether their robust supported values separate. The conditional null must be located once the prevalence search is included, and its movement with the width of $\PSet$ recorded. Sparse-positive designs should be examined for their effect on $C_{\mathrm{prev}}$ and $C_{\mathrm{rep}}$. A real paired application should then show whether a conclusion drawn at one reference prevalence reverses on the tie-averaged scale, and Section~\ref{sec:worst-case} identifies the case in which the prevalence set cannot change anything.

\section{Conclusion}

Average precision describes the purity of ranked selections in a target population. Its meaning changes with prevalence even when class-conditional ranking behavior remains fixed. Prior standardization evaluates that behavior throughout a declared set of target prevalences, and chance normalization gives every prevalence a common scale.

The prevalence-robust effect
\[
\inf_{\pi\in\PSet}\CNAP_\pi(D)
\]
records the level attained throughout the target set, and the replication-survival operator gives
\[
S_{K,\PSet}^{\mathrm{AP}}(\Regime)
=
\E_{\Regime}
\left[
\min_{1\leq j\leq K}
\inf_{\pi\in\PSet}
\CNAP_\pi(D_j)
\right].
\]
The supported value averages the maximal claim that survives every prevalence in every replication.

The three declarations of Section~\ref{sec:motivation} fix what that claim means. The prevalence set names the populations in which it must hold, the regime names the conditions under which it must hold, and the replication count says how many independent attempts at defeat it must survive. None of the three is estimated from the data, and a supported value quoted without them states nothing in particular.

The gain is uneven. When every contributing threshold of a single model lies above the ROC diagonal its profile is nondecreasing, and its worst case falls at the lower endpoint of an interval set. Paired differences can remain nonmonotone in the same evaluations, so the robust comparison retains content where the robust single-model summary reduces to a point-prevalence report.

For one fixed model evaluated on held-out outcomes, stratified resampling estimates the replication distribution and conditional permutation measures the chance level of the complete finite-design procedure. The calibrated value locates the model against that reference, and compatible marginal reports support separate design-relative interpretation when joint predictions are unavailable.

When both models are evaluated on the same units, the replication criterion is applied to
\[
\inf_{\pi\in\PSet}
\left\{
\CNAPta_{A,\pi}(D)-\CNAPta_{B,\pi}(D)
\right\},
\]
computed under the convention that averages over the orderings within each tie block. A replacement claim requires the supported advantage to exceed a declared margin, with a proposed outer lower bootstrap bound above that margin as the inferential criterion.

The two tasks end differently, and for a reason that is not a matter of taste. Both are threatened by a finite evaluation design whose chance level does not sit where the population scale says zero is. For an isolated model no convention moves that level to zero, so it has to be estimated and the reported value is a ratio against it. For a comparison the offending quantity is a difference of two such levels, and the tie-averaged convention makes that difference identically zero under the reference law. The comparison therefore keeps the raw scale, carries no positive population anchor attributable to the evaluation design under the declared reference law, and declines to choose when the replications disagree by more than the advantage in question. Prevalence profiles and the two costs explain these numbers when they are low, and belong with the diagnostics rather than in the claim.

What a study reports is one sentence:

\begin{quote}
This level survived every target prevalence declared relevant and every replication challenge represented by the regime.
\end{quote}

The value of the framework is that the sentence has a determinate meaning, and that a reader who disagrees with the declared populations or the declared regime can say so precisely.

\begin{thebibliography}{99}

\bibitem{bestgen2015random}
Bestgen, Y. (2015).
Exact expected average precision of the random baseline for system evaluation.
\emph{The Prague Bulletin of Mathematical Linguistics}, 103, 131--138.

\bibitem{bestgen2021fisher}
Bestgen, Y. (2021).
Improving measures of accuracy for detecting errors in texts.
\emph{Journal of Quantitative Linguistics}, 28, 293--310.

\bibitem{bareinboim2013transportability}
Bareinboim, E., \& Pearl, J. (2013).
A general algorithm for deciding transportability of experimental results.
\emph{Journal of Causal Inference}, 1(1), 107--134.

\bibitem{bouckaert2004estimating}
Bouckaert, R. R. (2004).
Estimating replicability of classifier learning experiments.
In \emph{Proceedings of the Twenty-First International Conference on Machine Learning}, 15.

\bibitem{cohen1960coefficient}
Cohen, J. (1960).
A coefficient of agreement for nominal scales.
\emph{Educational and Psychological Measurement}, 20, 37--46.

\bibitem{degeest2016online}
De Geest, R., Gavves, E., Ghodrati, A., Li, Z., Snoek, C., and Tuytelaars, T. (2016).
Online action detection.
In \emph{European Conference on Computer Vision}, 269--284.

\bibitem{elkan2001foundations}
Elkan, C. (2001).
The foundations of cost-sensitive learning.
In \emph{Proceedings of the Seventeenth International Joint Conference on Artificial Intelligence}, 973--978.

\bibitem{flach2015prg}
Flach, P. A. and Kull, M. (2015).
Precision-recall-gain curves: PR analysis done right.
In \emph{Advances in Neural Information Processing Systems}, 28.

\bibitem{hoeffding1948ustat}
Hoeffding, W. (1948).
A class of statistics with asymptotically normal distribution.
\emph{Annals of Mathematical Statistics}, 19, 293--325.

\bibitem{hosking1990lmoments}
Hosking, J. R. M. (1990).
L-moments: analysis and estimation of distributions using linear combinations of order statistics.
\emph{Journal of the Royal Statistical Society, Series B}, 52, 105--124.

\bibitem{lipton2018detecting}
Lipton, Z. C., Wang, Y.-X., \& Smola, A. (2018).
Detecting and correcting for label shift with black box predictors.
In \emph{Proceedings of the 35th International Conference on Machine Learning}, PMLR 80, 3122--3130.

\bibitem{mayo1996error}
Mayo, D. G. (1996).
\emph{Error and the Growth of Experimental Knowledge}.
University of Chicago Press.

\bibitem{nadeau2003inference}
Nadeau, C. and Bengio, Y. (2003).
Inference for the generalization error.
\emph{Machine Learning}, 52, 239--281.

\bibitem{ojala2010permutation}
Ojala, M. and Garriga, G. C. (2010).
Permutation tests for studying classifier performance.
\emph{Journal of Machine Learning Research}, 11, 1833--1863.

\bibitem{pearl2011transportability}
Pearl, J., \& Bareinboim, E. (2011).
Transportability of causal and statistical relations: A formal approach.
In \emph{Proceedings of the Twenty-Fifth AAAI Conference on Artificial Intelligence}, 247--254.

\bibitem{perdomo2020performative}
Perdomo, J. C., Zrnic, T., Mendler-D\"unner, C., \& Hardt, M. (2020).
Performative prediction.
In \emph{Proceedings of the 37th International Conference on Machine Learning}, PMLR 119, 7599--7609.

\bibitem{popper1959logic}
Popper, K. R. (1959).
\emph{The Logic of Scientific Discovery}.
Hutchinson.

\bibitem{popper1963conjectures}
Popper, K. R. (1963).
\emph{Conjectures and Refutations: The Growth of Scientific Knowledge}.
Routledge.

\bibitem{saerens2002adjusting}
Saerens, M., Latinne, P., \& Decaestecker, C. (2002).
Adjusting the outputs of a classifier to new a priori probabilities: A simple procedure.
\emph{Neural Computation}, 14(1), 21--41.

\bibitem{schmeidler1986integral}
Schmeidler, D. (1986).
Integral representation without additivity.
\emph{Proceedings of the American Mathematical Society}, 97(2), 255--261.

\bibitem{siblini2020calibration}
Siblini, W., Fr\'ery, J., He-Guelton, L., Obl\'e, F., and Wang, Y. (2020).
Master your metrics with calibration.
In \emph{Advances in Intelligent Data Analysis XVIII}, 457--469.

\bibitem{storkey2009training}
Storkey, A. (2009).
When training and test sets are different: Characterizing learning transfer.
In \emph{Dataset Shift in Machine Learning}, 3--28. MIT Press.

\bibitem{suyuan2015relationship}
Su, W., Yuan, Y., and Zhu, M. (2015).
A relationship between the average precision and the area under the ROC curve.
In \emph{Proceedings of the 2015 International Conference on The Theory of Information Retrieval}, 349--352.

\bibitem{wang1996premium}
Wang, S. (1996).
Premium calculation by transforming the layer premium density.
\emph{ASTIN Bulletin}, 26, 71--92.

\bibitem{yaari1987dual}
Yaari, M. E. (1987).
The dual theory of choice under risk.
\emph{Econometrica}, 55(1), 95--115.

\end{thebibliography}

\appendix

\section{Prior-standardized average precision}
\label{app:ap-definition}

\subsection{Empirical definition}

Let
\[
D=\{(s_i,y_i)\}_{i=1}^{n},
\qquad
y_i\in\{0,1\},
\]
with $n_+=\sum_i y_i>0$ and $n_-=n-n_+>0$. Let
\[
c_1>c_2>\cdots>c_H
\]
be the distinct observed score values. Threshold $h$ selects every observation with $s_i\geq c_h$. Define
\[
\operatorname{TP}_h
=
\sum_{i=1}^{n}
\mathbf{1}\{y_i=1,\ s_i\geq c_h\},
\qquad
\operatorname{FP}_h
=
\sum_{i=1}^{n}
\mathbf{1}\{y_i=0,\ s_i\geq c_h\},
\]
and
\[
r_h=\frac{\operatorname{TP}_h}{n_+},
\qquad
f_h=\frac{\operatorname{FP}_h}{n_-}.
\]
Prior-standardized precision at threshold $h$ is
\[
p_{\pi,h}
=
\frac{\pi r_h}
{\pi r_h+(1-\pi)f_h}.
\]
Set $r_0=0$. The non-interpolated prior-standardized AP is
\begin{equation}
\AP_\pi(D)
=
\sum_{h=1}^{H}
(r_h-r_{h-1})p_{\pi,h}.
\label{eq:empirical-ap}
\end{equation}
Equal scores enter as one threshold block. A threshold that adds only negative cases has zero recall increment. Its false positives affect later precision values.

\subsection{The tie-averaged convention}
\label{app:tie-convention}

The block convention above is what makes the finite-design displacement of Section~\ref{sec:calibration} depend on the score vector, and it is therefore what generates the differential anchor of Section~\ref{sec:paired-anchor}. Section~\ref{sec:paired-convention} requires a different convention for paired comparison. Let $\mathcal{O}(D)$ be the set of rankings of the evaluation units consistent with the observed scores, so that a vector of distinct scores admits one ranking and a vector with tie blocks of sizes $m_1,\ldots,m_J$ admits $\prod_j m_j!$ of them. Define
\begin{equation}
\APta_\pi(D)
=
\frac{1}{|\mathcal{O}(D)|}
\sum_{o\in\mathcal{O}(D)}
\AP_\pi(D_o),
\label{eq:tie-averaged-ap}
\end{equation}
where $D_o$ is the evaluation with the tie block resolved into the strict ranking $o$, and let $\CNAPta_\pi$ follow through Equation~\eqref{eq:cnap}.

Equation~\eqref{eq:jitter} is the same object in another form. Adding $\varepsilon U$ with $U$ continuous and $\varepsilon$ small enough to reorder only within blocks selects an element of $\mathcal{O}(D)$ uniformly at random, so the expectation over $U$ is the average in Equation~\eqref{eq:tie-averaged-ap}.

\paragraph{Proof of Proposition~\ref{prop:refinement}.}
Condition on $S_B$ and $Y$. Since $A$ agrees with $B$ on the order of the blocks, $\Rank_A$ takes values in $\mathcal{O}(D_B)$, and $\APta_{A,\pi}$ is by construction its conditional expectation of $\AP_\pi$ given $S_A$. By the tower property, $\E[\APta_{A,\pi}\mid S_B,Y]=\E[\AP_\pi(D_{\Rank_A})\mid S_B,Y]$, and the hypothesis makes $\Rank_A$ uniform on $\mathcal{O}(D_B)$, so this is Equation~\eqref{eq:tie-averaged-ap} evaluated at $B$. $\square$

The outcomes are held fixed throughout, so nothing here asks that the models be uninformative or that the units be exchangeable. Refinement neutrality is a statement about the convention and not about the design, and it constrains an expectation rather than any realized comparison.

\paragraph{Proof of Proposition~\ref{prop:uniformity}.}
Fix $n_+$, $n_-$ and $\pi$, and let the labels follow the reference law $P_0$ of Definition~\ref{def:label-condition}. For any single strict ranking $o$, the conditional law of the label sequence read along $o$ is uniform over the $\binom{n}{n_+}$ arrangements and is therefore the same for every $o$. Equation~\eqref{eq:empirical-ap} is a function of that sequence alone, and hence
\[
\E_0\left[\AP_\pi(D_o)\right]
=
\alpha(n_+,n_-,\pi)
\]
for every $o$, with $\alpha$ not depending on the ranking. Equation~\eqref{eq:tie-averaged-ap} is a convex combination of these terms, so $\E_0[\APta_\pi]=\alpha(n_+,n_-,\pi)$ whatever the tie structure, and the two models therefore share it. $\square$

\paragraph{Scope of Proposition~\ref{prop:uniformity}.}
The hypothesis cannot be weakened to invariance under the permutations a restricted design permits. Take four units in two matched pairs, each pair containing exactly one positive, so that the permitted permutations are the independent label swaps within pairs. Write $a_1,a_2$ for the first pair and $b_1,b_2$ for the second, let $\PSet=\{1/2\}$, and consider the two strict rankings
\[
o_1=(a_1,a_2,b_1,b_2),
\qquad
o_2=(a_1,b_1,a_2,b_2),
\]
for which tie averaging changes nothing. Averaging Equation~\eqref{eq:empirical-ap} over the four permitted assignments gives
\[
\E\left[\AP_{1/2}(o_1)\right]=\frac23,
\qquad
\E\left[\AP_{1/2}(o_2)\right]=\frac{11}{16},
\qquad
\E\left[\Dta_{1/2}\right]=\frac{1}{24}
\]
for the comparison of $o_2$ against $o_1$. Neither ranking carries any association with the outcomes and neither contains a tie, so the design anchor is not common to them and Proposition~\ref{prop:uniformity} fails. Proposition~\ref{prop:refinement} is untouched, since the two rankings are not related by refinement. Proposition~\ref{prop:conservative} nonetheless survives this pair, because the replication operator continues to charge for disagreement: the paired difference takes the values $1/3$, $0$, $0$ and $-1/6$ over the four assignments, so $S_2(\RobDta;\Regime_0)=-5/96<0$. The counterexample of Section~\ref{sec:paired-scope}, in which the tie blocks of one model coincide with the randomization strata, is the one that defeats both. That pair gives $179/2508\approx0.071$ at $\pi=1/10$, and extending the construction to $P$ matched pairs at $\pi=1/2$ gives $0.036$ at $n=8$, above $0.032$ at $n=16$, $0.013$ at $n=100$ and $0.008$ at $n=200$.

\paragraph{The obstruction described exactly.}
The obstruction is the group the design randomizes over. Let $G$ be the set of label permutations the design permits, so that the outcome vector satisfies $Y\circ g\overset{d}{=}Y$ for every $g\in G$. Two score vectors share a design anchor whenever one ranking is obtained from the other by relabeling units through some $g\in G$, since the law of the label sequence read along the two rankings is then the same. Definition~\ref{def:label-condition} is the case in which $G$ is the whole symmetric group, every ranking lies in one orbit, and Proposition~\ref{prop:uniformity} follows. For a smaller $G$ the property is checkable but is not met by two arbitrary models, which is why the method adopts the restriction of Section~\ref{sec:paired-scope} rather than a generalization. A design-specific convention would have to average over an orbit of $G$ in place of $\mathcal{O}(D)$, and whether a convention as clean as Equation~\eqref{eq:tie-averaged-ap} exists for a given design is left open here.


The value $\alpha$ is not $\pi$, so the finite-design displacement of Section~\ref{sec:calibration} is not removed by this convention. What is removed is its dependence on the score vector, which is the only part a paired difference sees. In the four-unit example of Section~\ref{sec:paired-anchor} both models return $\alpha=49/72$ under this convention, in place of $49/72$ and $1/2$ under the block convention.

\paragraph{Proof of Proposition~\ref{prop:conservative}.}
Two models evaluated on the same units share $n_+$, $n_-$ and $\pi$. Write $E$ for everything $\Regime_0$ draws before the outcomes, so that Definition~\ref{def:label-condition} holds conditionally on $E$. That proposition gives $\E_{\Regime_0}[\CNAPta_{A,\pi}\mid E]=\E_{\Regime_0}[\CNAPta_{B,\pi}\mid E]$ and hence $\E_{\Regime_0}[\Dta_\pi\mid E]=0$; averaging over $E$ gives $\E_{\Regime_0}[\Dta_\pi]=0$ for every $\pi\in\PSet$. By Equation~\eqref{eq:prevalence-selection} applied to $\Dta$,
\[
\E_{\Regime_0}[\RobDta]
\leq
\inf_{\pi\in\PSet}\E_{\Regime_0}[\Dta_\pi]
=0 ,
\]
and $S_K(T;\Regime)\leq\E_{\Regime}[T]$ for any $T$ because a minimum of $K$ copies is at most the first. Composing the two gives the statement. Replacing $\Dta$ by $-\Dta$ leaves $\E_{\Regime_0}[-\Dta_\pi]=0$ and the argument unchanged, which gives the reversed comparison. $\square$

In the four-unit example the tie-averaged value of $S_2$ is $-49/216$ in the orientation that the block convention reported as $+29/216$.

\paragraph{Costs of the convention.}
Three, and they should be weighed openly. Statistical power is lost, since Proposition~\ref{prop:conservative} bounds the null value below zero rather than at it, and Section~\ref{sec:orientation} quantifies the resulting no-verdict region. A fully tied score vector no longer evaluates to $\CNAP=0$, because $\alpha>\pi$; the population reading of zero in Equation~\eqref{eq:cnap} is unaffected, but the coincidence by which one particular finite score vector realized it is gone, and Section~\ref{sec:calibration} was already the reason not to rely on that coincidence. And Equation~\eqref{eq:tie-averaged-ap} is written as an enumeration over $\prod_j m_j!$ rankings, which is impractical for any but the smallest blocks; Appendix~\ref{app:tie-computation} gives an exact finite sum that evaluates it in $O(\sum_j m_j^2)$ operations, and that is what an implementation should use.

The convention applies to the paired quantity. Sections~\ref{sec:common-scale} through \ref{sec:calibration} retain the block convention of Equation~\eqref{eq:empirical-ap}, where the conditional permutation null of Section~\ref{sec:calibration} already locates whatever displacement the realized score vector produces. Because the two conventions differ, a paired advantage is not the difference of two reported single-model values, which Section~\ref{sec:paired} gives independent reasons not to compute. Both conventions must be named among the prespecified items of Section~\ref{sec:reporting}.

Using $H$ for the threshold count avoids conflict with the replication severity parameter $K$.

\subsection{Blockwise computation of tie-averaged AP}
\label{app:tie-computation}

Equation~\eqref{eq:tie-averaged-ap} is defined as an average over $\prod_j m_j!$ rankings and is never evaluated that way. Index the tie blocks $j=1,\ldots,J$ from the highest score downward, with $a_j$ positive units, $b_j$ negative units and $m_j=a_j+b_j$, and write
\[
A_j=\sum_{l<j}a_l,
\qquad
B_j=\sum_{l<j}b_l
\]
for the counts entering block $j$. Then
\begin{equation}
\APta_\pi(D)
=
\sum_{j=1}^{J}
\frac{a_j}{n_+}
\cdot
\frac{1}{m_j}
\sum_{k=1}^{m_j}
\ \sum_{t=\max(0,\,k-1-b_j)}^{\min(k-1,\,a_j-1)}
\frac{\binom{a_j-1}{t}\binom{b_j}{k-1-t}}{\binom{m_j-1}{k-1}}
\cdot
\frac{\pi r_{jkt}}
{\pi r_{jkt}+(1-\pi)f_{jkt}},
\label{eq:tie-averaged-closed}
\end{equation}
where
\[
r_{jkt}=\frac{A_j+t+1}{n_+},
\qquad
f_{jkt}=\frac{B_j+k-1-t}{n_-}.
\]

\paragraph{Proof.}
Resolving the score vector into a strict ranking makes every unit its own threshold, so the recall increment in Equation~\eqref{eq:empirical-ap} is $1/n_+$ at each positive unit and zero at each negative one, and
\[
\AP_\pi(D_o)
=
\frac{1}{n_+}
\sum_{i:\,y_i=1}
p_{\pi,\,\mathrm{rank}_o(i)} ,
\]
the prior-standardized precision at the position unit $i$ occupies in $o$. Averaging over $\mathcal{O}(D)$ acts independently within blocks, and the counts $A_j$ and $B_j$ entering block $j$ do not depend on how the earlier blocks were ordered internally. Fix a positive unit of block $j$. Under a uniformly random ordering of that block its within-block position $k$ is uniform on $\{1,\ldots,m_j\}$, and given $k$ the number $t$ of positives among the $k-1$ units of the block placed above it is hypergeometric on the remaining $a_j-1$ positives and $b_j$ negatives. Its position in the full ranking then carries $A_j+t+1$ positives and $B_j+k-1-t$ negatives at or above it, which gives $r_{jkt}$ and $f_{jkt}$. Summing over the $a_j$ interchangeable positive units of the block and over the blocks gives Equation~\eqref{eq:tie-averaged-closed}. $\square$

The inner double sum has $O(m_j^2)$ terms, so one evaluation costs $O(\sum_j m_j^2)$ operations, bounded by $O(n^2)$ for a fully tied vector and linear in $n$ when the blocks are short. A vector with no ties has $m_j=1$ throughout, the inner sums collapse, and Equation~\eqref{eq:tie-averaged-closed} reduces to Equation~\eqref{eq:empirical-ap}. Only the final factor depends on $\pi$, so the block counts and hypergeometric weights are computed once and reused across the prevalence search of Appendix~\ref{app:computation}. The identity has been verified against direct enumeration of Equation~\eqref{eq:tie-averaged-ap} on random block structures at several prevalences, to floating-point precision. The paired quantity requires two such evaluations per prevalence, one for each model, and no permutation layer.

\subsection{Weighted implementation}

The same quantity can be calculated by weighting the observed classes to represent prevalence $\pi$. Assign each positive observation weight
\[
w_+=\frac{\pi}{n_+}
\]
and each negative observation weight
\[
w_-=\frac{1-\pi}{n_-}.
\]
At threshold $h$, weighted precision is
\[
\frac{w_+\operatorname{TP}_h}
{w_+\operatorname{TP}_h+w_-\operatorname{FP}_h}
=
\frac{\pi r_h}
{\pi r_h+(1-\pi)f_h}.
\]
Weighted recall remains $r_h$, so stepwise weighted AP equals Equation~\eqref{eq:empirical-ap}.

\subsection{Behavior across target prevalences}

Combining Equations~\eqref{eq:cnap} and \eqref{eq:empirical-ap} gives
\begin{equation}
\CNAP_\pi(D)
=
\sum_{h=1}^{H}
(r_h-r_{h-1})
\frac{\pi(r_h-f_h)}
{\pi r_h+(1-\pi)f_h}.
\label{eq:empirical-cnap-profile}
\end{equation}
The derivative of one threshold contribution is
\begin{equation}
\frac{\partial}{\partial\pi}
\left\{
\frac{\pi(r_h-f_h)}
{\pi r_h+(1-\pi)f_h}
\right\}
=
\frac{f_h(r_h-f_h)}
{\{\pi r_h+(1-\pi)f_h\}^2}.
\label{eq:profile-derivative}
\end{equation}
A threshold with $r_h\geq f_h$ has a nondecreasing contribution. If this holds at every threshold with a positive recall increment, the complete CNAP profile is nondecreasing and its minimum over $[\pi_L,\pi_U]$ occurs at $\pi_L$. Profiles can have interior extrema when contributing thresholds lie on different sides of the ROC diagonal.

Two nondecreasing marginal profiles can have a nonmonotone difference. Prevalence-robust paired comparison can therefore remain informative when each single-model infimum occurs at the lower endpoint.

\section{Related precision standardizations}
\label{app:related}

Equation~\eqref{eq:precision-odds} isolates prevalence in the prior-odds factor. Prior standardization assigns that factor a declared value and retains the selected-set interpretation of precision. Precision-recall-gain analysis removes the factor from its threshold-level precision coordinate \cite{flach2015prg}.

Let
\[
p
=
\frac{\pi r}{\pi r+(1-\pi)f}.
\]
Precision gain is
\[
\operatorname{precG}
=
\frac{p-\pi}{(1-\pi)p}.
\]
Writing $D=\pi r+(1-\pi)f$ gives
\[
p-\pi
=
\frac{\pi(1-\pi)(r-f)}{D},
\qquad
(1-\pi)p
=
\frac{(1-\pi)\pi r}{D},
\]
and hence
\begin{equation}
\operatorname{precG}
=
\frac{r-f}{r}
=
1-\frac{f}{r}.
\label{eq:precision-gain}
\end{equation}
The prevalence terms cancel at the threshold level.

For $r>0$, the CNAP threshold contribution in Equation~\eqref{eq:threshold-cnap} satisfies
\begin{equation}
\lim_{\pi\to1}
\frac{\pi(r-f)}
{\pi r+(1-\pi)f}
=
1-\frac{f}{r}.
\label{eq:precision-gain-limit}
\end{equation}
Precision gain is therefore the high-prevalence limit of the threshold-level normalized precision used by CNAP.

The integral summaries remain different. CNAP averages its threshold contribution against recall increments. A precision-recall-gain area uses recall gain and its associated integration measure. The two constructions can consequently order models differently. The present framework retains target prevalence because selected-set purity in the populations represented by $\PSet$ carries the scientific interpretation of interest.

\section{Properties of prevalence-robust support}
\label{app:proofs}

\subsection{Joint-survival representation}

For a random variable $T\in[L,U]$,
\[
T
=
L+\int_L^U\mathbf{1}\{T>s\}\,ds.
\]
For independent copies $T_1,\ldots,T_K$,
\[
\min_jT_j
=
L+
\int_L^U
\prod_{j=1}^{K}
\mathbf{1}\{T_j>s\}\,ds.
\]
Taking expectations and using independence gives Equation~\eqref{eq:support-survival}.

Set $T=\inf_{\pi\in\PSet}\CNAP_\pi(D)$. The event $T>s$ is equivalent to
\[
\CNAP_\pi(D)>s
\quad
\text{for every }\pi\in\PSet.
\]
Substitution gives Equation~\eqref{eq:robust-joint-survival}.

\subsection{Challenge monotonicity}

Let $\Pi_1\subseteq\Pi_2$ and $K_1\leq K_2$. The identity $T=L+\int_L^U\mathbf{1}\{T>s\}\,ds$ holds for any $L$ below $T$, so both challenge spaces may be represented through Equation~\eqref{eq:refutation-integral} using the common lower bound $L=L_{\Pi_2}$, which lies below $L_{\Pi_1}$ by set monotonicity of the range. Couple the two by taking one realization of $K_2$ independent replications and letting the first $K_1$ of them serve the smaller challenge space. Then
\[
\Pi_1\times\{1,\ldots,K_1\}
\subseteq
\Pi_2\times\{1,\ldots,K_2\},
\]
so at every level $s$ the refutation class of Equation~\eqref{eq:refutation-class} satisfies $\RefSet_s^{(1)}\subseteq\RefSet_s^{(2)}$. Hence $\RefSet_s^{(2)}=\emptyset$ implies $\RefSet_s^{(1)}=\emptyset$, and
\[
\Prb_{\Regime}\left(\RefSet_s^{(2)}=\emptyset\right)
\leq
\Prb_{\Regime}\left(\RefSet_s^{(1)}=\emptyset\right)
\]
for every $s$. Integrating over $s$ in Equation~\eqref{eq:refutation-integral} gives Equation~\eqref{eq:challenge-monotonicity}. Set monotonicity is the case $K_1=K_2$ and replication monotonicity the case $\Pi_1=\Pi_2$.

\subsection{Ordering relative to the supported profile}

Let $X_\pi^{(j)}=\CNAP_\pi(D_j)$. For every realized set of replications,
\[
\min_j\inf_{\pi\in\PSet}X_\pi^{(j)}
=
\inf_{\substack{\pi\in\PSet\\1\leq j\leq K}}
X_\pi^{(j)}
=
\inf_{\pi\in\PSet}\min_jX_\pi^{(j)}.
\]
The two orderings of the minimum agree pathwise, so no loss occurs at this step. Write
\[
Y_\pi=\min_{1\leq j\leq K}X_\pi^{(j)}
\]
and take expectations. Since $\inf_\pi Y_\pi\leq Y_{\pi'}$ for every $\pi'\in\PSet$,
\[
\E\left[\inf_{\pi\in\PSet}Y_\pi\right]
\leq
\inf_{\pi\in\PSet}\E[Y_\pi],
\]
which proves Equation~\eqref{eq:strong-weak-order}. The gap arises entirely at this step. It is the cost of taking the expectation after the prevalence search rather than before.

\subsection{The $K=2$ decomposition}

For real values $a$ and $b$,
\[
\min(a,b)
=
\frac{a+b-|a-b|}{2}.
\]
Apply this identity to two independent copies of $Z_{\PSet}$ and take expectations to obtain Equation~\eqref{eq:robust-k2}.

\subsection{Two costs of the prevalence challenge}
\label{app:penalty}

Equation~\eqref{eq:two-costs} is an identity at every $K$, obtained by adding and subtracting $\E_{\Regime}[Z_{\PSet}]$ in the definitions of the two costs. The prevalence-challenge cost $C_{\mathrm{prev}}$ is nonnegative by Equation~\eqref{eq:prevalence-selection}. The replication-disagreement cost is nonnegative because $\min_{j\leq K}Z_{\PSet,j}\leq Z_{\PSet,1}$ in every realization, and it is nondecreasing in $K$ because enlarging the set of replications can only lower their minimum.

Only at $K=2$ does the cost reduce to two moments. The identity $\min(a,b)=\tfrac12(a+b-|a-b|)$ gives $C_{\mathrm{rep}}=\tfrac12\E|Z_{\PSet,1}-Z_{\PSet,2}|$, the second L-moment $\lambda_2$ of Appendix~\ref{app:families}. For general $K$, write
\[
\alpha_r=\int_0^1Q(v)(1-v)^r\,dv,
\]
so that Equation~\eqref{eq:quantile-weight} below reads $\SRob=K\alpha_{K-1}$ and
\[
C_{\mathrm{rep}}=\lambda_1-K\alpha_{K-1}.
\]
Higher L-moments enter as $K$ grows. At $K=3$, for instance, $S_{3,\PSet}^{\mathrm{AP}}(\Regime)=\lambda_1-\tfrac32\lambda_2+\tfrac12\lambda_3$, so the disagreement between replications is no longer summarized by their mean absolute difference alone.

Estimation uses the same bootstrap array that produces $\SRobHat$. The anchor requires the mean CNAP profile across replications, evaluated on the shared prevalence grid, together with its minimizer. Note that this minimizer is a property of the averaged profile and generally differs from any individual $\pi^\dagger(D^{(b)})$. The two costs then follow by subtraction, since $\E[Z_{\PSet}]$ is the mean of the array and $\SRobHat$ is its U-statistic.

Separating them answers a question the supported value alone cannot. A result limited by $C_{\mathrm{prev}}$ is limited by an unstable worst-case population, which more evaluation units may not repair if the profile is genuinely flat near its minimum. A result limited by $C_{\mathrm{rep}}$ is limited by replication variability, which a larger evaluation usually does repair under a fixed-model regime. The two call for different responses.

\subsection{Affine equivariance}

For $a>0$ and constant $c$,
\[
\inf_{\pi\in\PSet}\{aX_\pi+c\}
=
a\inf_{\pi\in\PSet}X_\pi+c.
\]
The support operator also satisfies
\[
S_K(aT+c;\Regime)
=
aS_K(T;\Regime)+c.
\]
These identities explain why the common pointwise chance normalization is compatible with the prevalence infimum and replication support. A prevalence-dependent affine transformation would not commute with the infimum.

\subsection{Two established families}
\label{app:families}

Let $Q$ be the quantile function of $Z_{\PSet}$. The minimum of $K$ independent copies equals $Q(U_{(1)})$, where $U_{(1)}$ is the smallest of $K$ independent uniform variables and has density $K(1-v)^{K-1}$. Hence
\begin{equation}
S_{K,\PSet}^{\mathrm{AP}}(\Regime)
=
\int_0^1
Q(v)\,K(1-v)^{K-1}\,dv.
\label{eq:quantile-weight}
\end{equation}
The weight decreases in $v$, so the operator places its mass on the lower tail of the replication distribution.

Two identifications follow and supply properties without separate proof.

\paragraph{L-moments.}
The first two L-moments of a distribution \cite{hosking1990lmoments} are $\lambda_1=\E[Z_{\PSet}]$ and $\lambda_2=\tfrac12\E|Z_{\PSet,1}-Z_{\PSet,2}|$, the latter being half the Gini mean difference. Equation~\eqref{eq:robust-k2} is therefore the exact identity
\[
S_{2,\PSet}^{\mathrm{AP}}(\Regime)
=
\lambda_1-\lambda_2 .
\]
L-moment estimators are linear in order statistics, which is the origin of the sorted-sum form in Appendix~\ref{app:computation}, and they exist whenever the first moment does.

\paragraph{Distortion measures.}
Equation~\eqref{eq:support-survival} applies the distortion $u\mapsto u^{K}$ to the survival function of $Z_{\PSet}$, placing $S_K$ in the dual-power family of Wang \cite{wang1996premium}. The following properties hold.

\begin{itemize}
\item \emph{Monotonicity under first-order stochastic dominance.} If $Q_A(v)\geq Q_B(v)$ for all $v\in(0,1)$, then Equation~\eqref{eq:quantile-weight} gives $S_K(A;\Regime)\geq S_K(B;\Regime)$. Section~\ref{sec:paired} uses this property.
\item \emph{Monotonicity under the concave order.} The quantile weight is nonincreasing, so a mean-preserving spread of the replication distribution cannot raise the supported value.
\item \emph{Translation equivariance and positive homogeneity.} These restate the affine identities above.
\item \emph{Comonotone additivity.} If $X$ and $Y$ are comonotone, then $Q_{X+Y}=Q_X+Q_Y$ and $S_K(X+Y)=S_K(X)+S_K(Y)$.
\item \emph{Superadditivity.} For any $X$ and $Y$ evaluated on the same replications, $\min_j(X_j+Y_j)\geq\min_jX_j+\min_jY_j$, so
\[
S_K(X+Y;\Regime)
\geq
S_K(X;\Regime)+S_K(Y;\Regime).
\]
Taking $X=A_{\PSet}-B_{\PSet}$ and $Y=B_{\PSet}$ gives the first inequality of Equation~\eqref{eq:robust-marginal-order}.
\end{itemize}

These representations place the functional in familiar families and supply the properties used elsewhere in the paper. The next subsection shows that the replication-survival criterion selects it from within them.

\subsection{Why the distortion is a power}
\label{app:characterization}

Let $\rho$ assign a real number to the law of a replication effect $T$ taking values in $[L,U]$. Suppose $\rho$ is law invariant, monotone under first-order stochastic dominance, normalized so that $\rho(c)=c$ for a constant claim, and additive over comonotone pairs. The Schmeidler representation then gives
\begin{equation}
\rho(T)
=
L
+
\int_L^U
g\left(\Prb(T>s)\right)
ds
\label{eq:distortion-representation}
\end{equation}
for a unique nondecreasing $g:[0,1]\to[0,1]$ with $g(0)=0$ and $g(1)=1$ \cite{schmeidler1986integral,yaari1987dual}. The value $g(u)$ is the weight the summary places on a level that one replication retains with probability $u$.

Now impose joint survival. A level is retained by two independent replications only when each retains it, so if they retain it with probabilities $u_1$ and $u_2$ it is jointly retained with probability $u_1u_2$. Requiring the weight on the jointly retained level to be the product of the weights on the separately retained ones gives
\begin{equation}
g(u_1u_2)=g(u_1)g(u_2),
\qquad
u_1,u_2\in[0,1].
\label{eq:multiplicative-cauchy}
\end{equation}
Write $u=e^{-v}$ for $v\geq0$ and $h(v)=g(e^{-v})$. Equation~\eqref{eq:multiplicative-cauchy} becomes $h(v_1+v_2)=h(v_1)h(v_2)$ with $h$ nonincreasing and $h(0)=g(1)=1$, whose solutions other than the zero function are $h(v)=e^{-cv}$ with $c\geq0$. Hence $g(u)=u^{c}$ for some $c\geq0$, and the normalization $g(0)=0$ rules out $c=0$.

The exponent is not obtained from this equation, and the argument does not need it to be. Appendix~\ref{app:proofs} fixes it directly: $S_K$ is the expected minimum of $K$ independent replications, and a level that one replication retains with probability $u$ is retained by all $K$ with probability $u^{K}$, so $g(u)=u^{K}$ and $\rho=S_K$ returns Equation~\eqref{eq:support-survival}. What Equation~\eqref{eq:multiplicative-cauchy} adds is that the powers are the only distortions closed under joint survival, and its work is therefore to exclude the neighboring functionals rather than to identify $K$.

Two neighbors show what is excluded. Using the correspondence $\varphi(v)=g'(1-v)$ between a distortion and the quantile weight of Equation~\eqref{eq:quantile-weight}, the lower quantile at level $\alpha$ has distortion
\[
g(u)=\mathbf{1}\{u>1-\alpha\},
\]
and the expected shortfall averaging the worst $\alpha$ fraction of replications has
\[
g(u)=\frac{\max\{u-(1-\alpha),0\}}{\alpha}.
\]
Both satisfy the four conditions leading to Equation~\eqref{eq:distortion-representation}, and neither satisfies Equation~\eqref{eq:multiplicative-cauchy}. At $\alpha=1/2$ the first gives $g(0.8)g(0.6)=1$ against $g(0.48)=0$, and the second gives $g(0.8)^2=0.36$ against $g(0.64)=0.28$. Both are legitimate summaries of a replication distribution. Neither expresses a requirement that independent replications jointly retain a level, which is the question Section~\ref{sec:support} asks.

\subsection{Marginal and paired robust support}

Let $U_\pi$ and $V_\pi$ be any two prevalence-indexed profiles, with $U_{\PSet}=\inf_{\pi\in\PSet}U_\pi$, $V_{\PSet}=\inf_{\pi\in\PSet}V_\pi$ and $\underline{T}_{\PSet}=\inf_{\pi\in\PSet}\{U_\pi-V_\pi\}$. Choosing any prevalence minimizing $U_\pi$ gives
\[
\underline{T}_{\PSet}
\leq
U_{\PSet}-V_\pi
\leq
U_{\PSet}-V_{\PSet},
\]
so $U_{\PSet}-V_{\PSet}\geq\underline{T}_{\PSet}$ pointwise and monotonicity of $S_K$ gives $S_K(U_{\PSet}-V_{\PSet};\Regime)\geq S_K(\underline{T}_{\PSet};\Regime)$, the second inequality of Equation~\eqref{eq:robust-marginal-order}. The first is the superadditivity of Appendix~\ref{app:families} with $X=U_{\PSet}-V_{\PSet}$ and $Y=V_{\PSet}$.

\subsection{Why paired null centering fails}
\label{app:paired-null}

\paragraph{The four-unit computation.}
Let $n=4$ with $n_+=n_-=2$ and $\PSet=\{1/2\}$. At $\pi=1/2$ the weights of Appendix~\ref{app:ap-definition} are $w_+=w_-=1/4$, so Equation~\eqref{eq:empirical-ap} reduces to unweighted stepwise AP. Model $A$ has four distinct scores; writing $i_1<i_2$ for the ranks of the two positive cases,
\[
\AP_{1/2}(A)
=
\frac12
\left(
\frac{1}{i_1}+\frac{2}{i_2}
\right),
\]
which over the six equally likely rank pairs gives the values listed in Section~\ref{sec:paired-anchor}, with sum $49/12$ and mean $49/72$. Model $B$ assigns one score to every unit, giving a single threshold block with $r_1=f_1=1$, precision $p_{1/2,1}=1/2$ and $\AP_{1/2}(B)=1/2$ in every assignment. By Equation~\eqref{eq:cnap} at $\pi=1/2$, $\CNAP_{1/2}=2\AP_{1/2}-1$, so $\CNAP_{1/2}(B)=0$ identically and $\Delta_{1/2}$ takes the values
\[
1,\quad
\frac23,\quad
\frac12,\quad
\frac16,\quad
0,\quad
-\frac16 ,
\]
with mean $13/36$. For two independent replications, $\Prb(\min=v_{(i)})=(13-2i)/36$ for the $i$th smallest of six equally likely values, and $S_2(\RobDelta;\Regime_0)=29/216$.

Under the canonical fixed-model regime of Section~\ref{sec:estimation}, each assignment is instead followed by the stratified bootstrap, which draws two positives with replacement from the two positive units and two negatives from the two negative units. Enumerating all $6\times16$ combinations and applying $S_2$ to each assignment's bootstrap law gives $719/2304$. The two numbers answer the same question under two regimes and neither is the other's correction.

\paragraph{Larger cases.}
The four-unit example is extreme in its second model, and the effect is not confined to that extreme. Under the same construction with eight units and three positives, a model with eight distinct scores compared against one with four tie blocks of two gives $\E[\Delta_{1/2}]\approx0.088$ and $S_2\approx0.009$; reversing which model is coarser gives $\E[\Delta_{1/2}]\approx-0.090$ and $S_2\approx-0.175$; and at $\pi=0.1$ with two positives the first comparison gives $\E[\Delta_{0.1}]\approx0.100$ and $S_2\approx0.019$. Neither model in any of these pairs carries signal, and the second case shows that the replication-disagreement cost sometimes compounds the distortion rather than offsetting it.


\paragraph{Why no additive correction is available.}
It is natural to try to repair the block convention by measuring the displacement under a joint permutation and subtracting it. The attempt fails, and the way it fails identifies what is wrong.

Take four units with two positives, $\PSet=\{1/2\}$, and
\[
s_A=(4,3,2,1),
\qquad
s_B=(2,2,1,1),
\]
with the two highest-scored units positive. Both models separate the classes perfectly, so every stratified bootstrap replication has $\CNAP_A=\CNAP_B=1$ and $\Delta\equiv0$. The raw supported difference is exactly zero in both orientations, which is the correct verdict.

The joint-permutation anchor is not zero, and it is not symmetric. Exact enumeration of the nested procedure gives $35/384$ for $A-B$ and $-295/1152$ for $B-A$. Subtracting them would report $-35/384\approx-0.091$ in one orientation and $+295/1152\approx+0.256$ in the other, so the coarser model would be credited with a supported advantage over a model whose predictions are identical to its own on every replication.

The mechanism is visible in the decomposition of the anchor. Writing it through Equation~\eqref{eq:robust-k2}, the anchor is $\E_{\Regime_0}[\Delta]-C_{\mathrm{rep}}$ evaluated under the permutation. The first term is antisymmetric in the orientation, taking the values $\pm0.1736$ here; the second is a mean absolute difference and is therefore identical in both orientations, taking the value $0.0825$. Subtracting the anchor consequently adds $C_{\mathrm{rep}}$ back in \emph{both} directions. Since $C_{\mathrm{rep}}$ is the replication disagreement that Equation~\eqref{eq:support-operator} exists to charge for, the correction undoes the severity mechanism, and it does so symmetrically. Score vectors exist for which both centered orientations are positive at once: $s_A=(2,2,1,1)$ against $s_B=(2,1,1,1)$ with the two lowest-scored units positive gives $+0.125$ and $+0.153$.

The deeper reason is that the displacement is not an additive constant. At perfect separation $\CNAP=1$ is the ceiling and no displacement is possible, while under no association the displacement is at its largest. A correction measured at the second condition and applied at the first therefore removes something that is not present. This is why Section~\ref{sec:paired-convention} removes the term by definition rather than by estimation, and why the resulting criterion is conservative rather than calibrated.

\paragraph{Coherence and the no-verdict region.}
For any integrable $T$ and any $K$, $\min_{j\leq K}(-T_j)=-\max_{j\leq K}T_j$, so
\[
S_K(T;\Regime)+S_K(-T;\Regime)
=
\E_{\Regime}
\left[
\min_{j\leq K}T_j-\max_{j\leq K}T_j
\right]
\leq0,
\]
which proves Equation~\eqref{eq:orientation} and shows that at most one orientation can report a supported advantage. At $K=2$ the identity $\max(a,b)-\min(a,b)=|a-b|$ gives a no-verdict width of $\E|T_1-T_2|=2C_{\mathrm{rep}}(T)$. The centered quantity of the preceding paragraph satisfies no such bound, because the two orientations are shifted by anchors that are not negatives of one another.

\subsection{Relation to a standard-error adjustment}

Throughout this paper $\pi$ denotes a target prevalence. In this subsection alone the circle constant also appears, and it is written $\varpi$ to keep the two apart.

Suppose $Z_{\PSet}$ is normally distributed with mean $\mu$ and standard deviation $\sigma$. Two independent copies have difference
\[
Z_{\PSet,1}-Z_{\PSet,2}
\sim
N(0,2\sigma^2),
\]
so
\[
\frac12
\E
\left[
|Z_{\PSet,1}-Z_{\PSet,2}|
\right]
=
\frac{\sigma}{\sqrt{\varpi}},
\qquad
\varpi=3.14159\ldots
\]
Equation~\eqref{eq:robust-k2} becomes
\[
S_{2,\PSet}^{\mathrm{AP}}(\Regime)
=
\mu-\frac{\sigma}{\sqrt{\varpi}}.
\]
This numerical relation belongs to the normal family. The general replication-disagreement cost depends on the complete replication distribution and is not a universal multiple of its standard deviation.

\subsection{The support distribution and frequency-qualified levels}

Let
\[
M_K
=
\min\{Z_{\PSet,1},\ldots,Z_{\PSet,K}\}.
\]
Its distribution records the level jointly retained by $K$ replications. Prevalence-robust supported performance is $\E[M_K]$.

A frequency-qualified retained level can be defined by
\[
m_\gamma^{(K)}
=
\sup
\left\{
m:
\Prb(M_K\geq m)\geq\gamma
\right\}.
\]
If $Q_{\PSet}$ is the quantile function of $Z_{\PSet}$ and the distribution is continuous, then
\begin{equation}
m_\gamma^{(K)}
=
Q_{\PSet}\left(1-\gamma^{1/K}\right).
\label{eq:frequency-level}
\end{equation}
This quantity answers a different question from the mean retained level. It reports a level attained jointly with frequency $\gamma$ under the regime. It remains data-conditional when estimated from a bootstrap replication array and is not a confidence bound for a population parameter.

\subsection{Contraction of the replication distribution}

Suppose a sequence of evaluation designs makes $Z_{\PSet}$ converge in probability to a stable value $z_0$ and the effects remain uniformly bounded. Then
\[
S_{K,\PSet}^{\mathrm{AP}}(\Regime)
\longrightarrow
z_0.
\]
Larger evaluations can produce this contraction when the regime redraws independent units from stable class-conditional score distributions. Environmental variation, clustered dependence, and repeated model development can prevent it. Replication counts therefore belong in the regime description even though the supported effect contains no explicit sample-size multiplier.

\subsection{Prevalence-specific replacement margins}
\label{app:variable-margin}

Some applications assign different practical meaning to the same CNAP difference at different prevalences. A prevalence-specific margin $d(\pi)$ can replace the constant $d$ of Equation~\eqref{eq:replacement} through
\begin{equation}
S_K\left(
\inf_{\pi\in\PSet}
\{\Dta_\pi-d(\pi)\};
\Regime
\right)>0.
\label{eq:variable-margin}
\end{equation}
Since $d(\pi)$ is deterministic, Proposition~\ref{prop:conservative} applies to the shifted quantity unchanged. The margin function must be declared with the prevalence set, and population-level evidence then requires an outer lower bound above zero for Equation~\eqref{eq:variable-margin} in place of a bound above $d$.


\section{Regime construction}
\label{app:regime}

This appendix writes out the mechanisms of Equation~\eqref{eq:regime-graph} and shows what each regime of Table~\ref{tab:regimes} holds.

Let $j$ index replications and $i$ index evaluation units within one. The mechanisms are
\[
E_j\sim P_E,
\qquad
D_{\mathrm{train}}^{(j)}\mid E_j\sim P_{\mathrm{train}\mid E},
\qquad
\widehat f^{(j)}=\mathcal{A}\!\left(D_{\mathrm{train}}^{(j)}\right),
\]
\[
U_{ij}\mid E_j\sim P_{U\mid E},
\qquad
(X_{ij},Y_{ij})\mid U_{ij}\sim P_{X,Y\mid U},
\qquad
S_{ij}=\widehat f^{(j)}(X_{ij}),
\]
with $\mathcal{A}$ the complete development procedure, including fitting, tuning, and model selection. A fixed-model regime holds $E$ and $\widehat f$ and redraws $U$. A new-environment regime releases $E$ and redraws the units within each drawn environment, so cluster resampling at the environment preserves the dependence among its units. A repeated-development regime releases $D_{\mathrm{train}}$ and reruns $\mathcal{A}$, which is why it costs a complete refit in every replication and why no bootstrap of one realized score vector approximates it.

Two conditions on this construction are used in the body.

The target-prevalence set stays in the computed part of the challenge space as long as the invariance of Equation~\eqref{eq:selection-diagram} holds. Each generated dataset yields class-conditional threshold rates, and Equation~\eqref{eq:reference-precision} evaluates those rates throughout $\PSet$ without further sampling. Where it fails, $\pi$ stops being a free argument of the evaluation and becomes a property of $E_j$. The environment law $P_E$ must then place mass on environments whose prevalences span $\PSet$, and the analysis requires data from those environments. A regime that redraws units within one environment cannot recover this, whatever range of $\pi$ is declared.

The reference law $\Regime_0$ of Section~\ref{sec:reference-law} leaves every mechanism above as written and replaces only the assignment of outcome labels within each evaluation, drawing them as Definition~\ref{def:label-condition} specifies conditionally on everything already generated. This is the sense in which it alters one part of the regime and holds the rest, and it is what Proposition~\ref{prop:conservative} refers to.

\section{Efficient computation}
\label{app:computation}

\subsection{All-pairs support for $K=2$}

Sort the prevalence-robust replication effects:
\[
z_{(1)}\leq\cdots\leq z_{(B)}.
\]
The value $z_{(i)}$ is the minimum in each pair with a larger-indexed value, so
\begin{equation}
\widehat S_{2,\PSet,B}^{\,\mathrm{AP}}
=
\frac{2}{B(B-1)}
\sum_{i=1}^{B-1}
(B-i)z_{(i)}.
\label{eq:sorted-pairs}
\end{equation}
Sorting costs $O(B\log B)$ and the weighted sum costs $O(B)$.

\subsection{General $K$}

For the order-$K$ U-statistic, $z_{(i)}$ is the minimum for every subset containing it and choosing the remaining $K-1$ values from the $B-i$ larger values. Hence
\begin{equation}
\widehat S_{K,\PSet,B}^{\,\mathrm{AP}}
=
\binom{B}{K}^{-1}
\sum_{i=1}^{B-K+1}
\binom{B-i}{K-1}
z_{(i)}.
\label{eq:sorted-general-k}
\end{equation}

\subsection{Monte Carlo standard error}

For $K=2$, define the leave-one replication projection
\[
\widehat h_{1,-i}(z_i)
=
\frac{1}{B-1}
\sum_{j\ne i}
\min(z_i,z_j).
\]
A first-order Monte Carlo standard error is
\begin{equation}
\widehat{\operatorname{se}}_{\mathrm{MC}}
=
\frac{2}{\sqrt B}
\operatorname{sd}_{1\leq i\leq B}
\left\{
\widehat h_{1,-i}(z_i)
\right\}.
\label{eq:mcse}
\end{equation}
Prefix sums over the sorted array evaluate all projections in linear time after sorting.

\subsection{Prevalence evaluation}

For a finite set with $G$ values, evaluating every prevalence within every bootstrap replication costs $O(BGH)$ after threshold counts are available, where $H$ is the number of distinct score thresholds. The values share the same $r_h$ and $f_h$, so vectorization over $\pi$ avoids rebuilding the ranking.

For a continuous interval, Equation~\eqref{eq:empirical-cnap-profile} supplies a smooth one-dimensional objective away from irrelevant zero-recall thresholds. Endpoint values and stationary points can be checked directly. A numerical optimizer should be verified against a dense grid during implementation.

Conditional calibration repeats the robust support procedure for each permutation. The permutations are independent computational jobs. Monte Carlo stability of the null mean usually requires fewer permutations than accurate estimation of a small tail probability. The paired comparison of Section~\ref{sec:paired} has no permutation layer at all, and its cost is one bootstrap array per orientation plus the tie-averaging of Equation~\eqref{eq:tie-averaged-ap}, which should use Equation~\eqref{eq:tie-averaged-closed} rather than enumeration.

\section{Calibration and interval considerations}
\label{app:calibration}

\subsection{Perfect-separation fixed point}

If every positive score exceeds every negative score, every bootstrap sample preserves that ordering. At each target prevalence,
\[
\AP_\pi=1,
\qquad
\CNAP_\pi=1.
\]
The prevalence infimum and every replication minimum equal one. The canonical raw and calibrated supported values therefore attain their upper anchor.

\subsection{Conditional validity}

The fixed-prediction permutation distribution is valid when the evaluation outcomes satisfy the design's own label symmetry with respect to the fixed scores under the null and the permutation respects that symmetry. The procedure conditions on the realized score vector, class counts, and tie structure. It does not account for outcome information used during model development.

The robust null includes the search over $\PSet$. Every permutation must use the same prevalence set, numerical optimization, resampling counts, and support severity as the observed calculation. Altering the set or optimization after inspecting the observed profile breaks the correspondence.

\subsection{Outer intervals}

An outer bootstrap sample yields a new empirical distribution of class-conditional scores. The complete inner estimator produces one value of the target statistic for that outer sample. Percentile intervals are simple. Basic, studentized, and bias-corrected intervals may respond differently to skewness and changes in the minimizing prevalence.

For paired superiority, the outer resampling preserves the pairing between model scores. No null is recomputed within the outer layer, since Section~\ref{sec:paired-convention} estimates no correction. A lower interval endpoint above the margin is the evidence the replacement criterion asks for. It is offered here as a proposed procedure whose coverage remains to be established, not as an interval known to attain its nominal level: the target contains an infimum over prevalence with a possibly nonunique or moving minimizer, an inner bootstrap functional and an outer resampling layer, and percentile methods are not guaranteed at flat or multiple minima. What Proposition~\ref{prop:conservative} makes conservative is the reference behavior of the criterion, and it says nothing about the coverage of any interval around it. Simulation should examine coverage when the prevalence minimum occurs at an endpoint, in the interior, at several points, or changes under small perturbations, and a validity argument under a unique well-separated minimizer would leave the nonregular case open.

\section{Reporting templates}
\label{app:report-template}

\subsection{Single-model performance}

A concise report can follow this sequence:
\begin{quote}
The model was evaluated over the prespecified target-prevalence set $\PSet$. Its observed prevalence-robust CNAP was [value], with a minimizing prevalence of [value or set]. Under the declared [same-design or prospective] regime with replication counts [counts] and $K=[K]$, prevalence-robust supported CNAP was [value]. The conditional permutation null was [value], giving calibrated robust support [value]. The permutation $p$-value was [value] and the outer [level]\% interval was [limits].
\end{quote}
The accompanying methods identify the AP convention, resampling unit, optimization method, bootstrap and permutation counts, Monte Carlo errors, and separation of model development from evaluation.

\subsection{Paired replacement decision}

A paired report can follow this sequence:
\begin{quote}
Model $A$ was compared with model $B$ on the same evaluation units over $\PSet$. The observed worst-prevalence CNAP advantage was [value], attained at [prevalence]. Under the declared regime and $K=[K]$, supported superiority was [value] in favor of $A$ and [value] in the reversed orientation, so the comparison [did or did not] reach a verdict. The prespecified replacement margin was [value], and the outer [level]\% interval was [limits].
\end{quote}
The report states whether the proposed outer lower bound exceeds the margin, notes that its coverage is not yet established, and describes the additional considerations entering the decision.

\end{document}
